\documentclass{article}
\usepackage[main,preprint]{neurips_2026}
\usepackage[utf8]{inputenc}
\usepackage[T1]{fontenc}
\usepackage{hyperref,url,booktabs,amsmath,amssymb,amsthm,graphicx,microtype,xcolor}
\newtheorem{proposition}{Proposition}
\newtheorem{lemma}{Lemma}
\newcommand{\D}{\Delta^{K-1}}
\newcommand{\R}{\mathbb R}
\newcommand{\Pp}{\mathbb P}
\newcommand{\E}{\mathbb E}
\newcommand{\ind}{\mathbf 1}
\newcommand{\clr}{\operatorname{clr}}
\newcommand{\DenseH}{$0.899\,{\scriptstyle\pm 0.003}$}
\newcommand{\DenseCoord}{$0.079\,{\scriptstyle\pm 0.002}$}
\newcommand{\DenseBonf}{$0.918\,{\scriptstyle\pm 0.003}$}
\newcommand{\HeteroHOne}{$0.588\,{\scriptstyle\pm 0.010}$}
\newcommand{\HeteroGroupOne}{$0.900\,{\scriptstyle\pm 0.007}$}
\newcommand{\HeteroOracleOne}{$0.675\,{\scriptstyle\pm 0.008}$}
\newcommand{\ZerosHBoundary}{$0.343\,{\scriptstyle\pm 0.017}$}
\newcommand{\ZerosHZBoundary}{$0.905\,{\scriptstyle\pm 0.007}$}
\newcommand{\ZerosClrVolume}{$0.941\,{\scriptstyle\pm 0.004}$}
\newcommand{\ZerosSmoothVolume}{$0.612\,{\scriptstyle\pm 0.015}$}
\newcommand{\MixedHBoundary}{$0.187\,{\scriptstyle\pm 0.011}$}
\newcommand{\MixedHZBoundary}{$0.910\,{\scriptstyle\pm 0.008}$}
\newcommand{\MixedHGZBoundary}{$0.921\,{\scriptstyle\pm 0.007}$}
\newcommand{\MixedGroupBoundary}{$0.585\,{\scriptstyle\pm 0.015}$}
\newcommand{\MixedSmoothBoundary}{$0.054\,{\scriptstyle\pm 0.009}$}
\newcommand{\MixedSmoothBigBoundary}{$0.134\,{\scriptstyle\pm 0.009}$}
\newcommand{\MixedHGZSize}{$0.490\,{\scriptstyle\pm 0.009}$}
\newcommand{\MixedGroupSize}{$0.178\,{\scriptstyle\pm 0.010}$}
\newcommand{\MixedOracleSize}{$0.038\,{\scriptstyle\pm 0.001}$}
\newcommand{\MixedHGZOne}{$0.911\,{\scriptstyle\pm 0.005}$}
\newcommand{\MixedHScaleOne}{$0.761\,{\scriptstyle\pm 0.010}$}
\newcommand{\ShiftH}{$0.764\,{\scriptstyle\pm 0.007}$}
\newcommand{\ShiftWeighted}{$0.900\,{\scriptstyle\pm 0.005}$}
\newcommand{\ShiftGroup}{$0.903\,{\scriptstyle\pm 0.004}$}
\newcommand{\ShiftWeightedOne}{$0.851\,{\scriptstyle\pm 0.007}$}
\newcommand{\ContrastA}{$-0.068\pm 0.002$}
\newcommand{\ContrastB}{$-0.631\pm 0.009$}
\newcommand{\ContrastC}{$0.733\pm 0.012$}
\newcommand{\ContrastD}{$0.295\pm 0.011$}
\newcommand{\ContrastE}{$0.151\pm 0.006$}
\newcommand{\ContrastF}{$-0.019\pm 0.002$}
\newcommand{\ContrastG}{$0.312\pm 0.013$}
\newcommand{\MixedCollapse}{34/100 (95\% Wilson interval [0.255, 0.437])}
\newcommand{\ShiftCollapse}{38/100 (95\% Wilson interval [0.291, 0.478])}
\newcommand{\ZerosCollapse}{100/100 (95\% Wilson interval [0.963, 1.000])}
\newcommand{\MainSeconds}{15.8}
\newcommand{\PeakRSS}{163.7}
\newcommand{\SparseInfinite}{71}
\newcommand{\SparseMean}{7.25}

\newcommand{\MaxIntegrationMeanDiff}{0.0008}
\newcommand{\MaxIntegrationSingleDiff}{0.0032}
\newcommand{\MaxAnalyticDiff}{0.00018}
\newcommand{\CheckSeconds}{43.6}

\title{When Marginal Coverage Hides the Boundary:\\Evaluating Conformal Regions for Compositional Responses}
\author{Autonomous research study}
\begin{document}
\maketitle

\begin{abstract}
Conformal regions for a composition must cover a whole vector, including its possible zero components. We study how evaluations based on marginal coverage and interior volume can obscure failures on the boundary of the simplex. We give an exact finite-calibration formula for the whole-simplex transition of an unsmoothed centered log-ratio score extended to infinity at zeros, and evaluate distance, density, coordinatewise, scaled, weighted, and stratified conformal constructions in controlled Dirichlet experiments. Across independent repetitions, pooled methods generally attain nominal marginal coverage under exchangeability but can severely under-cover high-dispersion groups and structural-zero responses. In one mixed three-component setting, pooled square-root distance regions cover boundary responses at \MixedHBoundary; a candidate-response and covariate stratification control reaches \MixedHGZBoundary, with substantially larger regions. Smoothing restores finite log-ratio scores but does not repair boundary coverage. A true-density oracle is efficient under a specified reference measure yet can also sacrifice a minority group. We recommend reporting coverage by relevant covariate and support strata together with stratum-specific size, and show exactly which guarantees each calibration choice provides. Results are synthetic and concern observed compositions, not latent probabilities inferred from single class labels.
\end{abstract}

\section{Introduction}
Probability vectors, topic proportions, and mixture fractions live on a closed simplex. A point estimate in this space does not specify a useful uncertainty region. Conformal prediction can attach a set with finite-sample marginal coverage, but that statement leaves substantial freedom in the shape, support, and allocation of errors across examples. A method can cover most compositions while systematically missing those with a zero component, and still satisfy its advertised guarantee.

We distinguish two problems that share a probability-vector predictor. In classification, the future response is a label and the conformal output is a subset of labels. Here the future response is an \emph{observed composition} $Y\in\D=\{y\in[0,1]^K:\sum_jy_j=1\}$. The output is a subset of this continuum and the coverage event is simultaneous inclusion of all coordinates. The softmax vector of a classifier is not itself a ground-truth composition. Without probability-valued targets or further observation assumptions, label coverage cannot establish coverage of a latent probability vector.

The familiar distinction between marginal and conditional coverage becomes especially consequential when the response distribution charges different faces of the simplex. Interior volume assigns every proper face size zero, regardless of its probability mass. Log-ratio scores create an additional difficulty: they are undefined at zeros. Common repairs include smoothing, but the impact of that repair on set size and support-specific reliability must be evaluated on the original target scale.

We ask four concrete questions. (Q1) At matched nominal marginal coverage, how do geometric scores differ in efficiency? (Q2) What exactly happens to raw log-ratio calibration as zero prevalence crosses the allowed error rate? (Q3) Which observable-group and candidate-response stratifications repair which coverage failures, and at what size cost? (Q4) Which conclusions survive a controlled change in group prevalence? Our hypotheses and original design were recorded in \texttt{PROTOCOL.md} before the main experiment. Two response-stratified controls were added after inspecting the initial results; we identify that analysis as exploratory.

Our contribution is a focused evaluation and an exact analysis of a boundary failure mechanism, not a new conformal algorithm. We provide a reproducible experiment with common predictions and paired randomness, an interior/face size diagnostic, and an oracle baseline under an explicit reference measure. We find that the best choice depends on the target guarantee and loss: marginal size optimization, group reliability, and boundary reliability are different objectives.

\section{Related work and scope}
Split conformal inference separates fitting from calibration and permits arbitrary fixed score functions. Its marginal guarantee and its difference from full conformal inference are established in \citet{lei2016}. \citet{angelopoulos2021} explicitly recommend checking adaptivity and subgroup coverage and describe scale normalization and group-balanced calibration. We use these constructions as baselines, rather than claiming that the need for conditional diagnostics is new. Full conformal and resampling methods principally address data reuse and refitting; comparing that tradeoff is outside this experiment, which holds the predictor fixed to isolate geometry and calibration.

Class-label sets such as RAPS \citep{angelopoulos2020} target a discrete response; their cardinality is not the volume of a region for a probability-valued response. More directly, \citet{javanmardi2024} construct conformal credal sets from probability-valued labels using distance and likelihood scores, investigate annotation noise, and evaluate set inefficiency. Thus conformalizing the simplex and using a density score are prior work. Our narrower distinction is between support-specific reliability and geometric size when exact zeros have positive probability, including a finite-calibration saturation formula and stratified efficiency measurements. Our regions need not be convex, so we use the term \emph{region} rather than imposing a convex credal-set interpretation.

Covariate shift breaks ordinary exchangeability. We implement the weighted split construction of \citet{tibshirani2019}, including the test-point atom at infinity, with known likelihood ratios. This is a controlled validity check, not an evaluation of estimated shift correction. To preserve independent research under the task's access constraint, we consulted only the five supplied starter papers cited here via direct PDF URLs; we made no thematic literature search. Our positioning is consequently relative to this restricted, explicitly documented reading set, not a comprehensive novelty claim.

\section{Regions, guarantees, and evaluation}
\label{sec:methods}
\paragraph{Split calibration.}
Conditional on a predictor $m(x)$ and any pilot estimates fitted independently of calibration, let $S_i=s(X_i,Y_i)$, $i=1,\ldots,n$. For miscoverage $\alpha$, write
\begin{equation}
 k=\lceil(n+1)(1-\alpha)\rceil,\quad
 \widehat q=S_{(k)},\quad C(x)=\{y\in\D:s(x,y)\le\widehat q\},
 \label{eq:split}
\end{equation}
where $S_{(n+1)}=+\infty$. We use nonrandomized inclusion of ties. Exchangeability gives $\Pp\{Y_{n+1}\in C(X_{n+1})\}\ge1-\alpha$, also for extended-real scores. Near-equality requires additional no-tie conditions; it need not hold with an atom at infinity. Coverage is averaged over calibration and test draws, not a guarantee for every realized calibration sample.

\paragraph{Geometric and coordinate scores.}
The main scores are $s_{\max}=\|y-m\|_\infty$, $s_2=\|y-m\|_2$, and $s_H=\|\sqrt y-\sqrt m\|_2$, where square roots are componentwise. The last differs from a conventional Hellinger distance only by a fixed factor, which cancels under calibration. For positive vectors define $\clr(y)_j=\log y_j-K^{-1}\sum_\ell\log y_\ell$. Raw CLR uses $s_{\rm CLR}=\|\clr(y)-\clr(m)\|_2$ in the interior and $+\infty$ otherwise. In this explicitly chosen extension, $\widehat q=+\infty$ yields the entire \emph{closed} simplex. This convention is necessary for validity of this construction; silently discarding boundary candidates would fail.

Smoothed CLR uses $\clr(y+\epsilon\boldsymbol1)$ on both candidates and centers, with $\epsilon=10^{-3}$ as the main setting and $10^{-4},10^{-2}$ as sensitivity controls. Dividing by $1+K\epsilon$ would cancel inside CLR. Membership always tests the original $y$; we do not replace the evaluation target by a smoothed vector. Every smoothed score is finite on the compact simplex, although its interpretation as a relative error changes with $\epsilon$ and $K$.

Coord calibrates $|Y_j-m_j(X)|$ separately at error $\alpha$ and intersects the resulting intervals with the simplex. It is a deliberately invalid control for \emph{joint} $1-\alpha$ coverage. Bonf uses error $\alpha/K$ for each coordinate and therefore has a union-bound joint guarantee. In both cases we use the exact order statistic, including infinity for undersized calibration sets.

\paragraph{Adapting calibration.}
H-scale divides $s_H$ by a pilot-estimated group RMS score $a_g$. H-group calibrates $s_H$ separately within a prespecified covariate group $G$. The first keeps a marginal guarantee but makes no group guarantee; the second guarantees coverage conditional on each $G$ under within-group exchangeability. H-weight uses the weighted empirical score distribution with masses proportional to $w(X_i)=dP_{\rm test,X}/dP_{\rm cal,X}(X_i)$ and an additional mass $w(x)$ at infinity \citep{tibshirani2019}.

The exploratory controls H-Z and H-GZ use the response partition $Z(y)=\ind\{\min_j y_j=0\}$. With calibration thresholds $q_z$ or $q_{g,z}$ they return
\begin{equation}
 C_Z(x)=\{y:s_H(x,y)\le q_{Z(y)}\},\qquad
 C_{GZ}(x)=\{y:s_H(x,y)\le q_{G(x),Z(y)}\}.
 \label{eq:response}
\end{equation}
No knowledge of the realized test target is used: the candidate determines its stratum during membership evaluation. This is the standard response-conditional conformal construction applied to a fixed support partition. It guarantees coverage within $Z$, or within each $(G,Z)$ cell, provided calibration and test examples are exchangeable within that cell. Small or empty cells can force infinite thresholds; this includes boundary strata absent from a finite calibration sample. Such a threshold includes only its candidate stratum, unlike the whole-simplex collapse of raw CLR. These partitions do not guarantee coverage within individual faces or at every $x$.

\paragraph{Reliability and size.}
We report marginal, $G$-conditional, and boundary/interior-conditional coverage. We also measure size relative to a declared reference distribution, not by raw score radius. For $K=3$, let $U_2$ be the uniform distribution on the triangle and $U_1$ an equally weighted mixture of uniform distributions on its three edges. Define
\begin{equation}
 V_r(C)=\E_{X_{\rm test}}\E_{U_r}[\ind\{U_r\in C(X_{\rm test})\}],\quad
 V_\lambda=\lambda V_2+(1-\lambda)V_1.
 \label{eq:size}
\end{equation}
Both are dimensionless fractions. $V_1$ does not replace actual boundary coverage: it weights edges uniformly whereas the target law need not. We present $(V_2,V_1)$ separately and use $V_{1/2}$ only when a scalar is needed. The weight is an evaluation preference, not a probability estimated from the data.

\begin{proposition}[Volume blindness]
For any measurable region $C$, removing all its boundary points leaves its $(K-1)$-dimensional volume unchanged but reduces coverage by $\Pp\{Y\in C(X)\cap\partial\D\}$. Thus interior volume alone cannot distinguish regions that differ in boundary reliability.
\label{prop:blind}
\end{proposition}
The statement follows because the boundary is a finite union of lower-dimensional faces and hence has zero interior volume. It motivates reporting coverage and size by support stratum; it does not select a universal scalar efficiency criterion.

\section{An exact boundary transition for raw CLR}
\label{sec:theory}
The next result quantifies why a valid implementation can be very conservative rather than merely numerically unstable.
\begin{proposition}[An atom at the maximum score]
\label{prop:atom}
Let calibration and test scores be i.i.d. on $\R\cup\{+\infty\}$ with $\Pp(S=+\infty)=\pi$ and an atomless conditional law given $S<+\infty$. For $1\le k\le n$ in Eq.~\eqref{eq:split},
\begin{align}
 \Pp(\widehat q=+\infty)&=\Pp\{\operatorname{Bin}(n,\pi)>n-k\},\label{eq:collapse}\\
 \Pp(S_{n+1}\le\widehat q)&=\frac{k}{n+1}+\left(1-\frac{k}{n+1}\right)
 \Pp\{\operatorname{Bin}(n+1,\pi)>n+1-k\}.\label{eq:atomcov}
\end{align}
For raw CLR with positive centers and scores finite on every interior target, the first event is exactly the whole-simplex event. Any finite-radius raw CLR region excludes every boundary target, so its population coverage cannot exceed $1-\pi$.
\end{proposition}
\paragraph{Proof sketch.}
The calibration quantile is infinite precisely when fewer than $k$ calibration scores are finite, giving Eq.~\eqref{eq:collapse}. Among $n+1$ scores, if at least $k$ are finite, averaging the leave-one-out inclusion event covers exactly the $k$ smallest finite scores. If fewer are finite, every leave-one-out threshold is infinite and all scores are covered. Exchangeability gives Eq.~\eqref{eq:atomcov}. A complete argument and the stratification guarantee are in Appendix~\ref{app:proofs}.

For fixed $\pi<\alpha$, collapse probability tends to zero as $n$ grows; for $\pi>\alpha$ it tends to one. At $\pi=\alpha\in(0,1)$ it tends to one half. More calibration data therefore sharpens this transition; it does not resolve it. When $\pi>\alpha$, obtaining a nontrivial region requires changing the score or its treatment of support. Smoothing is one option but supplies only a marginal guarantee. These statements concern the specified raw extension, not all log-ratio-based constructions.

\begin{figure}[t]
\centering\includegraphics[width=\linewidth]{../figures/transition.pdf}
\caption{Exact atom transition: lines are Eqs.~\eqref{eq:collapse}--\eqref{eq:atomcov}; points are 4,000 independent calibration arrays per setting. Left: 95\% Wilson intervals for collapse probability. Right: coverage is Rao--Blackwellized over the test score, with approximate 95\% Monte Carlo intervals. Finite scores are exponential, and infinite scores have probability $\pi$.}
\label{fig:transition}
\end{figure}

\section{Experimental design}
\label{sec:design}
\paragraph{Controlled compositions.}
We use entirely generated data, so exact support and noise strata are known. Independently draw $T\sim\operatorname{Unif}[-1,1]$ and $G\sim\operatorname{Bernoulli}(.25)$, set $X=(T,G)$, and define $m(T)$ through a fixed softmax (Appendix~\ref{app:dgp}). Draw $Y^*\mid X\sim\operatorname{Dirichlet}(\kappa_Gm(T))$. With probability $\pi_G$, choose one coordinate uniformly, set it to zero, and renormalize the remaining coordinates. Otherwise $Y=Y^*$. There are no vertices or multiple-zero targets. The same $m$ is used by every distance method: it is the conditional mean before deletion, and need not be the mean afterward. This intentionally isolates region construction from predictor fitting, at the expense of realism.

The four exchangeable settings use $(\kappa_0,\kappa_1;\pi_0,\pi_1)$ equal to Dense $(40,40;0,0)$, Heterogeneous $(160,8;0,0)$, Zeros $(40,40;.15,.15)$, and Mixed $(160,8;.02,.32)$. Mixed has marginal boundary probability $.095$, just below $\alpha=.1$, with boundary mass concentrated in the high-dispersion group. Shift uses the same conditional laws as Mixed but changes only test $P(G=1)$ to $.65$. Exact weights are $.65/.25$ for $G=1$ and $.35/.75$ otherwise. The label-conditional $Z$ law alone changes under this shift because the group mixture within each $Z$ changes.

\paragraph{Oracle density control.}
Oracle calibrates the negative log true conditional density relative to $\mu=\tfrac12 U_{K-1}+\tfrac12 U_{K-2}$, where the second distribution is uniform over codimension-one faces and uniform within each face. The deletion mechanism yields an analytic Dirichlet face density (Appendix~\ref{app:oracle}). Its population level sets minimize expected $\mu$-size subject to marginal coverage, illustrating an efficiency benchmark with full distribution knowledge. This is an established density-score construction, not a deployable estimated-density method or a promise of conditional coverage.

\paragraph{Replication and uncertainty.}
For every setting and $K\in\{3,8,20\}$ we run 100 independent repetitions, each with 1,000 pilot, 499 calibration, and 2,000 test examples. Pilot data estimate only the RMS scales; no threshold or hyperparameter uses test labels. SeedSequence streams derived from master seed 20261102 separate pilot, calibration, test, and integration randomness. Methods share all these draws within each repetition. Tables report means $\pm$ the half-width of a two-sided 95\% Student $t$ interval over the 100 repetition-level estimates; method contrasts use paired differences. These are pointwise Monte Carlo intervals for averages over new calibration samples, not simultaneous bands or guarantees for a particular fitted region. Full seed-level distributions are supplied.

For $K=3$ we estimate Eq.~\eqref{eq:size} using $2^{14}$ randomized Sobol joint samples over test covariates and uniform candidate points per repetition, sharing samples across methods. Higher-dimensional size is deliberately not estimated by unreliable rare-hit Monte Carlo. A $2^{18}$ check on 40 fixed cells changes mean volume estimates by at most \MaxIntegrationMeanDiff; an analytic box/simplex calculation independently checks Max (Appendix~\ref{app:checks}). Seven substantive unit checks pass. The final main experiment, including the atom sweep, took \MainSeconds{} seconds with four CPU workers on an Intel Xeon Gold 6230; per-process child peak RSS was \PeakRSS{} MiB. No GPU was used. All raw data, scripts, timings, and earlier runs are included.

\section{Results}
\label{sec:results}
\subsection{Geometry changes efficiency, not the basic rank guarantee}
\begin{table}[t]
\caption{Observed joint coverage and normalized size for $K=3$. Means $\pm$ 95\% interval half-width across 100 repetitions. $V_2$ measures interior area and $V_1$ measures edge length. H-Z and H-GZ are exploratory response-stratified controls. A displayed zero is rounded, not a claim of exact zero.}
\label{tab:geometry}\centering\small\setlength{\tabcolsep}{3pt}
\begin{tabular}{lcccccc}\toprule
 & \multicolumn{3}{c}{Dense} & \multicolumn{3}{c}{Zeros ($\pi=.15$)} \\
Method & Coverage & $V_2$ & $V_1$ & Coverage & $V_2$ & $V_1$ \\\midrule
Max & $0.899\,{\scriptstyle\pm 0.003}$ & $0.074\,{\scriptstyle\pm 0.001}$ & $0.051\,{\scriptstyle\pm 0.001}$ & $0.901\,{\scriptstyle\pm 0.003}$ & $0.149\,{\scriptstyle\pm 0.004}$ & $0.123\,{\scriptstyle\pm 0.004}$ \\
L2 & $0.898\,{\scriptstyle\pm 0.003}$ & $0.074\,{\scriptstyle\pm 0.001}$ & $0.053\,{\scriptstyle\pm 0.001}$ & $0.901\,{\scriptstyle\pm 0.003}$ & $0.151\,{\scriptstyle\pm 0.004}$ & $0.128\,{\scriptstyle\pm 0.005}$ \\
H & $0.899\,{\scriptstyle\pm 0.003}$ & $0.067\,{\scriptstyle\pm 0.001}$ & $0.000\,{\scriptstyle\pm 0.000}$ & $0.901\,{\scriptstyle\pm 0.003}$ & $0.309\,{\scriptstyle\pm 0.012}$ & $0.143\,{\scriptstyle\pm 0.009}$ \\
CLR & $0.902\,{\scriptstyle\pm 0.003}$ & $0.134\,{\scriptstyle\pm 0.003}$ & $0.000\,{\scriptstyle\pm 0.000}$ & $1.000\,{\scriptstyle\pm 0.000}$ & $1.000\,{\scriptstyle\pm 0.000}$ & $1.000\,{\scriptstyle\pm 0.000}$ \\
CLR$_{10^{-3}}$ & $0.902\,{\scriptstyle\pm 0.003}$ & $0.129\,{\scriptstyle\pm 0.002}$ & $0.000\,{\scriptstyle\pm 0.000}$ & $0.901\,{\scriptstyle\pm 0.003}$ & $0.941\,{\scriptstyle\pm 0.004}$ & $0.240\,{\scriptstyle\pm 0.010}$ \\
H-group & $0.901\,{\scriptstyle\pm 0.003}$ & $0.067\,{\scriptstyle\pm 0.001}$ & $0.000\,{\scriptstyle\pm 0.000}$ & $0.903\,{\scriptstyle\pm 0.003}$ & $0.317\,{\scriptstyle\pm 0.011}$ & $0.152\,{\scriptstyle\pm 0.009}$ \\
H-Z & $0.899\,{\scriptstyle\pm 0.003}$ & $0.067\,{\scriptstyle\pm 0.001}$ & $1.000\,{\scriptstyle\pm 0.000}$ & $0.902\,{\scriptstyle\pm 0.003}$ & $0.067\,{\scriptstyle\pm 0.001}$ & $0.744\,{\scriptstyle\pm 0.004}$ \\
H-GZ & $0.901\,{\scriptstyle\pm 0.003}$ & $0.067\,{\scriptstyle\pm 0.001}$ & $1.000\,{\scriptstyle\pm 0.000}$ & $0.903\,{\scriptstyle\pm 0.003}$ & $0.067\,{\scriptstyle\pm 0.001}$ & $0.750\,{\scriptstyle\pm 0.004}$ \\
Oracle & $0.898\,{\scriptstyle\pm 0.003}$ & $0.061\,{\scriptstyle\pm 0.001}$ & $0.000\,{\scriptstyle\pm 0.000}$ & $0.901\,{\scriptstyle\pm 0.003}$ & $0.102\,{\scriptstyle\pm 0.001}$ & $0.070\,{\scriptstyle\pm 0.003}$ \\
\bottomrule
\end{tabular}

\end{table}
Dense results illustrate the limited meaning of comparing coverage alone (Table~\ref{tab:geometry}). Geometric methods remain close to nominal. H attains coverage \DenseH{} and smaller interior area than raw CLR: the paired $V_2$(H)$-V_2$(CLR) contrast is \ContrastA. This is an empirical ranking under the chosen Dirichlet law and interior measure, not a universal ordering of metrics. For a homogeneous interior response, assigning size to unsupported edges penalizes some absolute-error regions as well.

Coordinatewise validity is insufficient: in Dense at $K=20$, Coord joint coverage is \DenseCoord{}, whereas Bonf achieves \DenseBonf. Dependence imposed by summing to one does not turn separately calibrated coordinate intervals into a jointly valid region. Bonf is a useful conservative baseline but generally cannot allocate error as efficiently as a single joint score. Appendix tables report every method, including both controls.

In Zeros, raw CLR collapses in \ZerosCollapse. Smoothed CLR avoids infinite thresholds but can be very large: $V_2$ is \ZerosClrVolume{} at $\epsilon=10^{-3}$, compared with \ZerosSmoothVolume{} at $10^{-2}$. The paired interior-volume difference between H and CLR$_{10^{-3}}$ is \ContrastB. These size statements must be read together with boundary reliability rather than treating all nominally valid regions as interchangeable.

\subsection{Boundary and group reliability are separate objectives}
In Heterogeneous, pooled H covers the high-dispersion group at \HeteroHOne; H-group reaches \HeteroGroupOne. Its paired $V_{1/2}$ size difference relative to H is \ContrastF: group balancing can improve conditional reliability while also reducing size by avoiding over-inflation for the majority. Oracle still covers $G=1$ at only \HeteroOracleOne. Even the true conditional density can allocate errors unevenly when the objective is marginal coverage and average size.

Zeros isolates the support issue without group dispersion. Pooled H's boundary coverage is \ZerosHBoundary, while H-Z reaches \ZerosHZBoundary. An approximately 90\% marginal region can spend much of its allowed error on the 15\% boundary subpopulation. Group balancing cannot repair this homogeneous within-group support failure because $G$ is unrelated to it.

\begin{figure}[t]
\centering\includegraphics[width=\linewidth]{../figures/conditional.pdf}
\caption{Mixed: marginal, high-dispersion-group, and boundary coverage differ markedly. Bars are pointwise 95\% intervals over repetitions. H-group targets $G$; H-Z targets zero status; H-GZ targets their intersection. Neither a density oracle nor smoothing promises these conditional targets.}
\label{fig:conditional}
\end{figure}

Mixed combines both sources of heterogeneity (Figure~\ref{fig:conditional} and Table~\ref{tab:reliability}). At $K=3$, H's boundary coverage is \MixedHBoundary, and H-group improves it only to \MixedGroupBoundary. H-scale's $G=1$ coverage, \MixedHScaleOne, shows that a pilot RMS correction need not normalize the score tails. H-Z improves boundary coverage to \MixedHZBoundary{} but can retain large group disparities; H-GZ reaches \MixedHGZBoundary{} on the boundary and \MixedHGZOne{} for $G=1$. The paired boundary-coverage gain of H-GZ relative to H is \ContrastC. These are expected consequences of their distinct target guarantees, not evidence that one construction dominates every other criterion.

The size cost can be substantial. Mixed $K=3$ gives $V_{1/2}=$\MixedHGZSize for H-GZ versus \MixedGroupSize{} for H-group (paired difference \ContrastG). The low-zero-probability $G=0,Z=1$ cell has mean calibration count \SparseMean; its threshold is infinite in \SparseInfinite{} of 100 repetitions. That correctly includes the entire boundary stratum for those queries. Pooling faces stabilizes boundary calibration but does not eliminate scarcity within the cross-partition.

Oracle attains $V_{1/2}=$\MixedOracleSize, but this small size is coupled to weak conditional coverage. Comparing it to H-GZ solely by size would compare different reliability objectives. A reference measure also matters: Appendix Figure~\ref{fig:weights} varies its interior/edge weight and retains the separate components in all raw outputs.

\begin{table}[t]
\caption{Reliability diagnostics at $K=8$. All entries are mean coverage $\pm$ 95\% interval half-width. H-weight equals H in the exchangeable Mixed setting; under Shift it uses exact density ratios. The oracle density is not weighted under shift.}
\label{tab:reliability}\centering\small\setlength{\tabcolsep}{3pt}
\begin{tabular}{lccccc}\toprule
 & \multicolumn{3}{c}{Mixed, $K=8$} & \multicolumn{2}{c}{Shift, $K=8$} \\
Method & Marginal & $G=1$ & Boundary & Marginal & $G=1$ \\\midrule
H & $0.898\,{\scriptstyle\pm 0.003}$ & $0.607\,{\scriptstyle\pm 0.012}$ & $0.413\,{\scriptstyle\pm 0.012}$ & $0.751\,{\scriptstyle\pm 0.007}$ & $0.619\,{\scriptstyle\pm 0.011}$ \\
CLR$_{10^{-3}}$ & $0.899\,{\scriptstyle\pm 0.003}$ & $0.620\,{\scriptstyle\pm 0.012}$ & $0.300\,{\scriptstyle\pm 0.012}$ & $0.756\,{\scriptstyle\pm 0.007}$ & $0.629\,{\scriptstyle\pm 0.010}$ \\
H-scale & $0.899\,{\scriptstyle\pm 0.003}$ & $0.842\,{\scriptstyle\pm 0.006}$ & $0.531\,{\scriptstyle\pm 0.010}$ & $0.870\,{\scriptstyle\pm 0.003}$ & $0.843\,{\scriptstyle\pm 0.005}$ \\
H-group & $0.901\,{\scriptstyle\pm 0.003}$ & $0.902\,{\scriptstyle\pm 0.006}$ & $0.624\,{\scriptstyle\pm 0.012}$ & $0.904\,{\scriptstyle\pm 0.003}$ & $0.905\,{\scriptstyle\pm 0.005}$ \\
H-Z & $0.901\,{\scriptstyle\pm 0.003}$ & $0.603\,{\scriptstyle\pm 0.011}$ & $0.908\,{\scriptstyle\pm 0.011}$ & $0.741\,{\scriptstyle\pm 0.007}$ & $0.601\,{\scriptstyle\pm 0.011}$ \\
H-GZ & $0.904\,{\scriptstyle\pm 0.003}$ & $0.906\,{\scriptstyle\pm 0.006}$ & $0.919\,{\scriptstyle\pm 0.010}$ & $0.906\,{\scriptstyle\pm 0.003}$ & $0.907\,{\scriptstyle\pm 0.005}$ \\
Oracle & $0.901\,{\scriptstyle\pm 0.003}$ & $0.649\,{\scriptstyle\pm 0.009}$ & $0.578\,{\scriptstyle\pm 0.012}$ & $0.761\,{\scriptstyle\pm 0.006}$ & $0.641\,{\scriptstyle\pm 0.008}$ \\
H-weight & $0.898\,{\scriptstyle\pm 0.003}$ & $0.607\,{\scriptstyle\pm 0.012}$ & $0.413\,{\scriptstyle\pm 0.012}$ & $0.903\,{\scriptstyle\pm 0.004}$ & $0.851\,{\scriptstyle\pm 0.006}$ \\
\bottomrule
\end{tabular}

\end{table}

\subsection{Raw and smoothed log ratios fail in different ways}
The independent score experiment in Figure~\ref{fig:transition} agrees with the exact binomial transition and coverage expressions. Raw CLR in Mixed $K=3$ collapses in \MixedCollapse, consistent with a zero probability close to the error budget. Both a finite radius that entirely excludes boundary compositions and an infinite radius that includes everything can occur across calibrations. Their average may look reassuring while neither output is attractive.

Smoothing changes this mechanism but supplies no support-conditional guarantee. Mixed $K=3$ boundary coverage at $\epsilon=10^{-3}$ is only \MixedSmoothBoundary; at $10^{-2}$ it is \MixedSmoothBigBoundary. Thus the finite-score repair should not be described as a repair of boundary reliability. Appendix Figure~\ref{fig:smoothing} reports the full declared sensitivity range, with uncertainty. We select no best $\epsilon$ after observing results. The exact zeros here are structural; data rounded to zero, censored measurements, or finite annotation counts would require a separate observation model.

\subsection{Shift correction repairs its specified target}
\begin{figure}[t]
\centering\includegraphics[width=\linewidth]{../figures/shift.pdf}
\caption{Changing test $P(G=1)$ from .25 to .65 breaks pooled marginal validity. Exact weighting repairs the marginal target; within-$G$ and within-$(G,Z)$ exchangeability preserve stratified validity. H-Z alone is not invariant to this shift. Intervals quantify variability over independent repetitions.}
\label{fig:shift}
\end{figure}
Under Shift, H's marginal coverage at $K=3$ falls to \ShiftH. Exact weighting restores it to \ShiftWeighted, and H-group obtains \ShiftGroup. At $K=8$ the paired improvement from weighting H is \ContrastE. Weighting does not equalize groups: $K=3$ H-weight covers $G=1$ at \ShiftWeightedOne. H-Z alone loses its boundary-conditional exchangeability when the group composition within $Z$ changes; H-GZ retains it because the conditional law given $(G,Z)$ is unchanged. These outcomes verify the assumptions behind the methods rather than suggesting robustness to arbitrary shift. In particular, an unweighted oracle density score does not retain marginal validity merely by knowing the true conditional response law.

\section{Discussion and limitations}
\label{sec:discussion}
Our results support a conditional recommendation. For an interior response with an adequate center, standard joint scores are valid and their efficiency should be compared under the geometry relevant to the decision. With exact zeros, inspect how the score defines membership on every relevant face and report support-specific coverage. Raw CLR with the specified infinity extension has a predictable saturation risk; smoothing offers finite scores but should be assessed as a different geometry. Use prespecified covariate or candidate-response stratification when the corresponding conditional coverage is required, budgeting calibration samples for the smallest cells. Under known covariate shift, correct weights address marginal validity; they do not promise conditional equality.

The evidence has deliberate limits. All datasets are synthetic and use a fixed analytic center; we do not establish performance for learned softmax models, topic models, annotation distributions, or empirical latent probabilities. We vary dimension and noise but only use one softmax family, a Dirichlet interior law, and single-coordinate deletion. We do not study vertices, multiple-zero patterns, continuous sparsity, dependence, estimation of shift weights, or model training instability. The fixed absolute smoothing parameter also has a different relative scale as $K$ increases. Size evidence is restricted to $K=3$; coverage at larger $K$ does not validate high-dimensional volume rankings. Face size is geometrical and may poorly represent task-specific error costs.

The boundary-stratified controls were added after the initial experiment; their comparisons are exploratory and do not constitute independent confirmation on a new benchmark. Confidence intervals are pointwise, capture Monte Carlo uncertainty under the declared mechanisms, and cannot justify universal method selection. The density oracle uses unavailable information and is a benchmark, not an implementable recommendation. Within a reliability objective, actual applications should select geometry and support partitions before evaluating on held-out targets, and should explicitly define whether the target is an observed composition or an unknown latent distribution.

\paragraph{Broader impact and research process.}
Support-specific diagnostics can expose harmful overconfidence when zero fractions correspond to consequential outcomes. Conversely, nominal marginal coverage, arbitrary size weights, or overly fine partitions could create a false sense of safety or excessive conservatism. Our artificial groups represent no protected population and do not establish fairness. No personal data or human participants were used. This study was autonomously designed, coded, analyzed, and drafted by an LLM-based research agent; no human review is represented. All reported experiments are executable CPU computations without model API calls, and the complete access and failure log is provided.

\clearpage
\bibliographystyle{plainnat}
\bibliography{references}

\clearpage
\appendix
\section{Proofs and precise targets}
\label{app:proofs}
\paragraph{Split rank lemma.}
Conditional on an independently fitted score, assign independent continuous tie-breaking marks to $n+1$ exchangeable scores on the extended real line. The rank of the test score in the resulting total order is uniform. Rank at most $k$ implies numerical inclusion under the nonrandomized threshold in Eq.~\eqref{eq:split}. This establishes coverage at least $k/(n+1)\ge1-\alpha$; for $k=n+1$ the set is the full candidate domain. Removing the tie-breaking marks can only increase inclusion. This standard argument applies to the arbitrary scores used here and does not assert an upper bound when numerical ties occur.

\paragraph{Proof of Proposition~\ref{prop:atom}.}
Let $B_n=\sum_{i=1}^n\ind\{S_i=+\infty\}$. There are $n-B_n$ finite calibration scores. Its $k$th order statistic is infinite if and only if $n-B_n<k$, which is $B_n>n-k$. Independence gives $B_n\sim\operatorname{Bin}(n,\pi)$, proving Eq.~\eqref{eq:collapse}.

For the coverage identity, let $N=n+1$ and condition on the unordered $N$-score sample. With probability one all finite scores are distinct. Write $F=N-B_N$ for their count. By exchangeability, the expected coverage equals the average of the $N$ inclusion events obtained by treating each score in turn as test and calibrating on the other $n$ scores. If $F\ge k$, a finite test score of rank $r\le k$ is covered: if there are enough remaining finite scores its rank lies below their $k$th score, and if not the threshold is infinite. A finite test score with $r>k$ is not covered because the remaining calibration sample contains $k$ strictly smaller finite scores. An infinite test score is not covered because the calibration still contains at least $k$ finite scores. Exactly $k$ test choices are therefore covered, giving average $k/N$.

If $F<k$, removing any score leaves fewer than $k$ finite calibration scores. Every leave-one-out threshold is infinite, so all $N$ choices are included, including the infinite scores under the specified extended-real comparison. Hence unconditional coverage is
\begin{equation}
 \frac{k}{N}\Pp(F\ge k)+\Pp(F<k)
 =\frac{k}{N}+\left(1-\frac{k}{N}\right)\Pp(B_N>N-k),
\end{equation}
which is Eq.~\eqref{eq:atomcov}. A finite raw CLR radius cannot include a boundary point; consequently the probability of inclusion for any fixed finite-radius rule is at most $\Pp(Y\notin\partial\D)=1-\pi$. This last claim does not bound a random rule that sometimes has infinite radius. The asymptotic transition follows from the law of large numbers for $\pi\ne\alpha$ and the central limit theorem with $k=\lceil N(1-\alpha)\rceil$ at equality.

\paragraph{Response-conditional calibration.}
Let $A(X,Y)$ be a fixed finite-valued partition, here $Z(Y)$ or $(G(X),Z(Y))$. Conditional on a cell's calibration count $n_a$ and on the test point being in cell $a$, its calibration pairs and the test pair are i.i.d. from the same cell-conditional distribution whenever that distribution is unchanged across calibration and testing. Applying the split rank lemma gives
\begin{equation}
 \Pp\{s_H(X,Y)\le q_a\mid A(X,Y)=a,n_a\}\ge1-\alpha,
\end{equation}
with the full stratum included if the rank exceeds $n_a$. Averaging over $n_a$ preserves the inequality. Region membership in Eq.~\eqref{eq:response} selects precisely that threshold for the true candidate, establishing the claim. This does not condition on a particular calibration draw or yield pointwise-$X$ coverage. For H-group use $A=G$; for H-GZ use the cross-partition. Under our shift, $Y\mid(T,G)$ and $T\mid G$ are unchanged. Therefore laws within $G$ and $(G,Z)$ are unchanged, whereas the law within $Z$ alone need not be.

\paragraph{Volume blindness.}
For fixed $x$, the indicator difference between $C(x)$ and $C(x)\setminus\partial\D$ is supported on the boundary, which has zero $(K-1)$-dimensional Lebesgue measure on the affine hyperplane. Their normalized interior sizes coincide. Their coverage indicator difference is $\ind\{Y\in C(X)\cap\partial\D\}$, giving Proposition~\ref{prop:blind} after expectation. The statement does not require any smoothness of $C$ or of the response distribution.

\section{Data generation, oracle density, and reproducibility}
\label{app:dgp}
For $j=0,\ldots,K-1$, define $\theta_j=2\pi j/K$ and
\begin{equation}
 \ell_j(t)=1.1\cos\theta_j+1.2t\sin\theta_j+.4t^2\cos(2\theta_j),\qquad
 m_j(t)=\frac{\exp\ell_j(t)}{\sum_r\exp\ell_r(t)}.
\end{equation}
The generator draws independent gamma variables with shapes $\kappa_G m_j(t)$ and unit scale and normalizes their sum. It then performs the disclosed zero/deletion step. Every run asserts nonnegativity, finiteness, and that exact zeros arise only from this step, preventing numerical gamma underflow from being mistaken for structural sparsity. Features and target draws are independent across pilot, calibration, and test sets. There is no external dataset or license dependency for the generated observations. Code is released under MIT and generated data under CC0; library versions and notices appear in the accompanying environment files.

H-scale uses $a_g=\sqrt{N_g^{-1}\sum_{i\in\mathrm{pilot},G_i=g}s_H(X_i,Y_i)^2}$. We do not tune this estimator, the concentrations, the smoothing grid, or the size weight using the test data. Calibration sets for every method are identical. The response strata use exact binary zero status, not an estimated near-zero tolerance. This distinction matters for real measurements.

\subsection{Density relative to a stratified measure}
\label{app:oracle}
Write $a_j=\kappa_g m_j(t)$, and let $\operatorname{Dir}_d(\cdot;a)$ denote the ordinary density in the first $d-1$ coordinates. Uniform normalized simplex density is $\Gamma(d)$ in the same coordinates. Dirichlet aggregation shows that deleting coordinate $j$ and normalizing leaves a Dirichlet vector with parameters $a_{-j}$, \emph{without} rescaling their total to $\kappa_g$. Each deleted coordinate has probability $\pi_g/K$.

For an interior point and a point on the relative interior of face $j$, respectively, the conditional density relative to $\mu=\tfrac12U_{K-1}+\tfrac12U_{K-2}$ is
\begin{align}
 f_\mu(y\mid t,g)&=2(1-\pi_g)\frac{\operatorname{Dir}_K(y;a)}{\Gamma(K)},
 &&y\in\operatorname{int}\D,\\
 f_\mu(y\mid t,g)&=2\pi_g\frac{\operatorname{Dir}_{K-1}(y_{-j};a_{-j})}{\Gamma(K-1)},
 &&y_j=0,\ y_{-j}>0.
\end{align}
The factors $1/K$ cancel between actual face probability and reference face-mixture weight. Lower faces carry zero probability under this DGP and this measure; the code sets their density to zero, an arbitrary harmless convention for the reported population metrics. Oracle calibrates $-\log f_\mu$. In no-zero settings it assigns infinite score to face candidates. Its behavior on faces is therefore well defined, unlike using an interior density formula alone at $y_j=0$.

\paragraph{Population efficiency target.}
Let $\nu=P_X\otimes\mu$ and $f=dP_{X,Y}/d\nu$. A level set $C_c=\{(x,y):f(y\mid x)\ge c\}$ has minimal $\nu$-size among measurable sets with at least its coverage, for $c>0$. Indeed, $(\ind_{C_c}-\ind_A)(f-c)\ge0$ pointwise, so
\begin{equation}
 P(C_c)-P(A)\ge c\{\nu(C_c)-\nu(A)\}.
\end{equation}
If $P(A)\ge P(C_c)$, this implies $\nu(A)\ge\nu(C_c)$. A possible density tie can be partially included to match a target. This is a population statement for a specified $P_X$, reference measure, and coverage, not an exact optimality assertion for randomly calibrated finite-sample regions. Ordinary calibration of this score still needs calibration/test exchangeability, explaining its failure in Shift. Conditional reliability would impose additional constraints and hence a different optimization problem.

\subsection{Uncertainty, seeds, and generated outputs}
For each metric let $M_r$ be its estimate in repetition $r$. The displayed half-width is $t_{99,.975}\operatorname{sd}(M_1,\ldots,M_{100})/\sqrt{100}$. Conditional coverages first divide by the number of test targets in that stratum within each repetition. Every nonempty response stratum has positive probability in the scenarios where it is reported; no boundary metric is defined in Dense or Heterogeneous. Size integration randomness is also independent between repetitions. Constant empirical metrics can give zero-width $t$ intervals; they should not be interpreted as a universal bound on an unobserved rare event. Collapse proportions in the text instead use Wilson binomial intervals. Paired contrasts are stored with unrounded interval limits in \texttt{results/paired\_contrasts.csv}. We do not perform multiple-testing-adjusted winner selection.

All seed streams are derived by NumPy SeedSequence from $[20261102,\mathrm{scenario\ index},K,r]$, with four children in pilot/calibration/test/integration order. Scenario order is Dense, Heterogeneous, Zeros, Mixed, Shift and $r=0,\ldots,99$. The atom sweep uses seed 20261103 and 4,000 actual score arrays per $(n,\pi)$, for $n\in\{99,499,1999\}$ and the complete probability grid in \texttt{src/run\_experiments.py}. Conditional test coverage for a realized threshold $q$ is analytically integrated as $(1-\pi)(1-e^{-q})$ for finite $q$ and one otherwise. This reduces test noise without replacing simulated calibration by the exact answer.

Run \texttt{bash reproduce.sh} from the artifact root after installing the pinned dependencies. This regenerates experiments, all figures and tables, numerical checks, and the compiled paper. \texttt{README.md} documents a results-only build, exact environment commands, expected runtime, and the script for every figure/table. The initial main run is retained in \texttt{results/initial}; the unchanged methods agree exactly with the final run on all common reported metrics. A separate two-repetition smoke run and its different method inventory are retained and are not included in paper estimates.

\section{Numerical validation and sensitivity}
\label{app:checks}
For 10 fixed repetitions of each of the four exchangeable $K=3$ settings, we repeat volume integration with $2^{18}$ Sobol points using the same scramble. The maximum absolute change in a mean over these 10 repetitions, across scenario/method/size-component combinations, is \MaxIntegrationMeanDiff; the maximum change in an individual estimate is \MaxIntegrationSingleDiff. These checks took \CheckSeconds{} seconds and are provided in \texttt{results/integration\_check.csv}. They validate numerical stability only for these cells, not a uniform integration error bound.

For an independent analytic check, a box $l_j\le y_j\le u_j$ intersected with the $K$-component simplex has normalized interior volume
\begin{equation}
 \sum_{A\subseteq\{1,\ldots,K\}}(-1)^{|A|}
 \left[1-\sum_jl_j-\sum_{j\in A}(u_j-l_j)\right]_+^{K-1},
\end{equation}
obtained by translation and inclusion--exclusion on the upper bounds. For Max set $l_j=\max(0,m_j-q)$ and $u_j=\min(1,m_j+q)$. Numerical quadrature over $T$ evaluates expected volume without sampling candidate compositions. For one calibration of each exchangeable setting, the largest difference from $2^{18}$-point Sobol integration is \MaxAnalyticDiff. This also checks the normalized volume convention. Unit checks independently verify the oracle density against SciPy's Dirichlet density and exhaustively check the atom coverage identity over leave-one-out choices in a small sample.

\begin{figure}[h]
\centering\includegraphics[width=\linewidth]{../figures/smoothing.pdf}
\caption{Declared smoothing ablation at $K=3$. Every point recalibrates its own score on the same calibration data. Larger smoothing changes both size components and can change boundary coverage while preserving only the marginal rank guarantee. All error bars are 95\% intervals over repetitions.}
\label{fig:smoothing}
\end{figure}
\begin{figure}[h]
\centering\includegraphics[width=\linewidth]{../figures/volume_weights.pdf}
\caption{Sensitivity to the reference-measure preference at $K=3$. Lines combine the actual measured $V_2,V_1$ means; they are not additional experiments. Each point is an explicit mixture preference, not an empirically identified universal notion of size. Uncertainty for the two components is in the raw summaries.}
\label{fig:weights}
\end{figure}

\clearpage
\section{Complete results}
\label{app:full}
All entries are means $\pm$ pointwise 95\% interval half-width over 100 independent repetitions. A dash denotes an undefined or intentionally unestimated metric. $V_{1/2}$ is estimated only for $K=3$. No boundary targets occur in Dense or Heterogeneous. H-Z/H-GZ are exploratory additions. Tables are generated directly from \texttt{results/runs.csv}; unrounded results, seed-level variability, and both size components are in the accompanying CSV files.

\subsection{Dense}
\begin{center}\small\setlength{\tabcolsep}{4pt}\begin{tabular}{lccccc}\toprule
Method & Marginal & $G=0$ & $G=1$ & Boundary & $V_{1/2}$ \\\midrule
\multicolumn{6}{l}{$K=3$} \\
Max & $0.899\,{\scriptstyle\pm 0.003}$ & $0.899\,{\scriptstyle\pm 0.003}$ & $0.898\,{\scriptstyle\pm 0.004}$ & -- & $0.063\,{\scriptstyle\pm 0.001}$ \\
L2 & $0.898\,{\scriptstyle\pm 0.003}$ & $0.898\,{\scriptstyle\pm 0.003}$ & $0.897\,{\scriptstyle\pm 0.004}$ & -- & $0.063\,{\scriptstyle\pm 0.001}$ \\
H & $0.899\,{\scriptstyle\pm 0.003}$ & $0.899\,{\scriptstyle\pm 0.003}$ & $0.899\,{\scriptstyle\pm 0.004}$ & -- & $0.033\,{\scriptstyle\pm 0.000}$ \\
CLR & $0.902\,{\scriptstyle\pm 0.003}$ & $0.902\,{\scriptstyle\pm 0.003}$ & $0.902\,{\scriptstyle\pm 0.004}$ & -- & $0.067\,{\scriptstyle\pm 0.001}$ \\
CLR$_{10^{-3}}$ & $0.902\,{\scriptstyle\pm 0.003}$ & $0.902\,{\scriptstyle\pm 0.003}$ & $0.902\,{\scriptstyle\pm 0.004}$ & -- & $0.065\,{\scriptstyle\pm 0.001}$ \\
H-scale & $0.899\,{\scriptstyle\pm 0.003}$ & $0.900\,{\scriptstyle\pm 0.003}$ & $0.897\,{\scriptstyle\pm 0.005}$ & -- & $0.033\,{\scriptstyle\pm 0.000}$ \\
H-group & $0.901\,{\scriptstyle\pm 0.003}$ & $0.901\,{\scriptstyle\pm 0.003}$ & $0.902\,{\scriptstyle\pm 0.006}$ & -- & $0.034\,{\scriptstyle\pm 0.000}$ \\
H-Z & $0.899\,{\scriptstyle\pm 0.003}$ & $0.899\,{\scriptstyle\pm 0.003}$ & $0.899\,{\scriptstyle\pm 0.004}$ & -- & $0.533\,{\scriptstyle\pm 0.000}$ \\
H-GZ & $0.901\,{\scriptstyle\pm 0.003}$ & $0.901\,{\scriptstyle\pm 0.003}$ & $0.902\,{\scriptstyle\pm 0.006}$ & -- & $0.534\,{\scriptstyle\pm 0.000}$ \\
Oracle & $0.898\,{\scriptstyle\pm 0.003}$ & $0.898\,{\scriptstyle\pm 0.003}$ & $0.898\,{\scriptstyle\pm 0.005}$ & -- & $0.030\,{\scriptstyle\pm 0.000}$ \\
Coord & $0.784\,{\scriptstyle\pm 0.004}$ & $0.784\,{\scriptstyle\pm 0.004}$ & $0.784\,{\scriptstyle\pm 0.006}$ & -- & $0.039\,{\scriptstyle\pm 0.001}$ \\
Bonf & $0.925\,{\scriptstyle\pm 0.003}$ & $0.925\,{\scriptstyle\pm 0.003}$ & $0.925\,{\scriptstyle\pm 0.003}$ & -- & $0.069\,{\scriptstyle\pm 0.001}$ \\
H-weight & $0.899\,{\scriptstyle\pm 0.003}$ & $0.899\,{\scriptstyle\pm 0.003}$ & $0.899\,{\scriptstyle\pm 0.004}$ & -- & $0.033\,{\scriptstyle\pm 0.000}$ \\
CLR$_{10^{-2}}$ & $0.902\,{\scriptstyle\pm 0.003}$ & $0.902\,{\scriptstyle\pm 0.003}$ & $0.902\,{\scriptstyle\pm 0.004}$ & -- & $0.050\,{\scriptstyle\pm 0.001}$ \\
CLR$_{10^{-4}}$ & $0.902\,{\scriptstyle\pm 0.003}$ & $0.902\,{\scriptstyle\pm 0.003}$ & $0.902\,{\scriptstyle\pm 0.004}$ & -- & $0.067\,{\scriptstyle\pm 0.001}$ \\
\midrule
\multicolumn{6}{l}{$K=8$} \\
Max & $0.897\,{\scriptstyle\pm 0.003}$ & $0.897\,{\scriptstyle\pm 0.003}$ & $0.899\,{\scriptstyle\pm 0.004}$ & -- & -- \\
L2 & $0.898\,{\scriptstyle\pm 0.003}$ & $0.897\,{\scriptstyle\pm 0.003}$ & $0.900\,{\scriptstyle\pm 0.004}$ & -- & -- \\
H & $0.901\,{\scriptstyle\pm 0.003}$ & $0.901\,{\scriptstyle\pm 0.003}$ & $0.901\,{\scriptstyle\pm 0.004}$ & -- & -- \\
CLR & $0.900\,{\scriptstyle\pm 0.003}$ & $0.900\,{\scriptstyle\pm 0.003}$ & $0.901\,{\scriptstyle\pm 0.004}$ & -- & -- \\
CLR$_{10^{-3}}$ & $0.900\,{\scriptstyle\pm 0.003}$ & $0.900\,{\scriptstyle\pm 0.003}$ & $0.902\,{\scriptstyle\pm 0.004}$ & -- & -- \\
H-scale & $0.901\,{\scriptstyle\pm 0.003}$ & $0.901\,{\scriptstyle\pm 0.003}$ & $0.899\,{\scriptstyle\pm 0.004}$ & -- & -- \\
H-group & $0.903\,{\scriptstyle\pm 0.003}$ & $0.902\,{\scriptstyle\pm 0.003}$ & $0.906\,{\scriptstyle\pm 0.006}$ & -- & -- \\
H-Z & $0.901\,{\scriptstyle\pm 0.003}$ & $0.901\,{\scriptstyle\pm 0.003}$ & $0.901\,{\scriptstyle\pm 0.004}$ & -- & -- \\
H-GZ & $0.903\,{\scriptstyle\pm 0.003}$ & $0.902\,{\scriptstyle\pm 0.003}$ & $0.906\,{\scriptstyle\pm 0.006}$ & -- & -- \\
Oracle & $0.901\,{\scriptstyle\pm 0.003}$ & $0.901\,{\scriptstyle\pm 0.003}$ & $0.901\,{\scriptstyle\pm 0.004}$ & -- & -- \\
Coord & $0.423\,{\scriptstyle\pm 0.005}$ & $0.423\,{\scriptstyle\pm 0.005}$ & $0.424\,{\scriptstyle\pm 0.006}$ & -- & -- \\
Bonf & $0.906\,{\scriptstyle\pm 0.003}$ & $0.906\,{\scriptstyle\pm 0.003}$ & $0.908\,{\scriptstyle\pm 0.004}$ & -- & -- \\
H-weight & $0.901\,{\scriptstyle\pm 0.003}$ & $0.901\,{\scriptstyle\pm 0.003}$ & $0.901\,{\scriptstyle\pm 0.004}$ & -- & -- \\
CLR$_{10^{-2}}$ & $0.900\,{\scriptstyle\pm 0.003}$ & $0.901\,{\scriptstyle\pm 0.003}$ & $0.900\,{\scriptstyle\pm 0.004}$ & -- & -- \\
CLR$_{10^{-4}}$ & $0.900\,{\scriptstyle\pm 0.003}$ & $0.900\,{\scriptstyle\pm 0.003}$ & $0.901\,{\scriptstyle\pm 0.004}$ & -- & -- \\
\midrule
\multicolumn{6}{l}{$K=20$} \\
Max & $0.896\,{\scriptstyle\pm 0.003}$ & $0.897\,{\scriptstyle\pm 0.003}$ & $0.896\,{\scriptstyle\pm 0.004}$ & -- & -- \\
L2 & $0.897\,{\scriptstyle\pm 0.003}$ & $0.898\,{\scriptstyle\pm 0.003}$ & $0.896\,{\scriptstyle\pm 0.003}$ & -- & -- \\
H & $0.898\,{\scriptstyle\pm 0.003}$ & $0.899\,{\scriptstyle\pm 0.003}$ & $0.896\,{\scriptstyle\pm 0.004}$ & -- & -- \\
CLR & $0.899\,{\scriptstyle\pm 0.003}$ & $0.900\,{\scriptstyle\pm 0.003}$ & $0.896\,{\scriptstyle\pm 0.004}$ & -- & -- \\
CLR$_{10^{-3}}$ & $0.900\,{\scriptstyle\pm 0.003}$ & $0.899\,{\scriptstyle\pm 0.003}$ & $0.901\,{\scriptstyle\pm 0.004}$ & -- & -- \\
H-scale & $0.898\,{\scriptstyle\pm 0.003}$ & $0.899\,{\scriptstyle\pm 0.003}$ & $0.896\,{\scriptstyle\pm 0.004}$ & -- & -- \\
H-group & $0.899\,{\scriptstyle\pm 0.003}$ & $0.899\,{\scriptstyle\pm 0.003}$ & $0.900\,{\scriptstyle\pm 0.006}$ & -- & -- \\
H-Z & $0.898\,{\scriptstyle\pm 0.003}$ & $0.899\,{\scriptstyle\pm 0.003}$ & $0.896\,{\scriptstyle\pm 0.004}$ & -- & -- \\
H-GZ & $0.899\,{\scriptstyle\pm 0.003}$ & $0.899\,{\scriptstyle\pm 0.003}$ & $0.900\,{\scriptstyle\pm 0.006}$ & -- & -- \\
Oracle & $0.900\,{\scriptstyle\pm 0.003}$ & $0.899\,{\scriptstyle\pm 0.003}$ & $0.903\,{\scriptstyle\pm 0.004}$ & -- & -- \\
Coord & $0.079\,{\scriptstyle\pm 0.002}$ & $0.079\,{\scriptstyle\pm 0.002}$ & $0.078\,{\scriptstyle\pm 0.003}$ & -- & -- \\
Bonf & $0.918\,{\scriptstyle\pm 0.003}$ & $0.919\,{\scriptstyle\pm 0.003}$ & $0.917\,{\scriptstyle\pm 0.004}$ & -- & -- \\
H-weight & $0.898\,{\scriptstyle\pm 0.003}$ & $0.899\,{\scriptstyle\pm 0.003}$ & $0.896\,{\scriptstyle\pm 0.004}$ & -- & -- \\
CLR$_{10^{-2}}$ & $0.896\,{\scriptstyle\pm 0.004}$ & $0.896\,{\scriptstyle\pm 0.004}$ & $0.896\,{\scriptstyle\pm 0.004}$ & -- & -- \\
CLR$_{10^{-4}}$ & $0.900\,{\scriptstyle\pm 0.003}$ & $0.901\,{\scriptstyle\pm 0.003}$ & $0.900\,{\scriptstyle\pm 0.004}$ & -- & -- \\
\bottomrule
\end{tabular}
\end{center}
\clearpage
\subsection{Heterogeneous}
\begin{center}\small\setlength{\tabcolsep}{4pt}\begin{tabular}{lccccc}\toprule
Method & Marginal & $G=0$ & $G=1$ & Boundary & $V_{1/2}$ \\\midrule
\multicolumn{6}{l}{$K=3$} \\
Max & $0.898\,{\scriptstyle\pm 0.003}$ & $1.000\,{\scriptstyle\pm 0.000}$ & $0.594\,{\scriptstyle\pm 0.010}$ & -- & $0.093\,{\scriptstyle\pm 0.002}$ \\
L2 & $0.898\,{\scriptstyle\pm 0.003}$ & $1.000\,{\scriptstyle\pm 0.000}$ & $0.594\,{\scriptstyle\pm 0.010}$ & -- & $0.093\,{\scriptstyle\pm 0.002}$ \\
H & $0.897\,{\scriptstyle\pm 0.003}$ & $1.000\,{\scriptstyle\pm 0.000}$ & $0.588\,{\scriptstyle\pm 0.010}$ & -- & $0.070\,{\scriptstyle\pm 0.002}$ \\
CLR & $0.896\,{\scriptstyle\pm 0.003}$ & $1.000\,{\scriptstyle\pm 0.000}$ & $0.586\,{\scriptstyle\pm 0.010}$ & -- & $0.130\,{\scriptstyle\pm 0.005}$ \\
CLR$_{10^{-3}}$ & $0.896\,{\scriptstyle\pm 0.003}$ & $1.000\,{\scriptstyle\pm 0.000}$ & $0.586\,{\scriptstyle\pm 0.010}$ & -- & $0.126\,{\scriptstyle\pm 0.004}$ \\
H-scale & $0.900\,{\scriptstyle\pm 0.003}$ & $0.895\,{\scriptstyle\pm 0.003}$ & $0.914\,{\scriptstyle\pm 0.005}$ & -- & $0.053\,{\scriptstyle\pm 0.001}$ \\
H-group & $0.901\,{\scriptstyle\pm 0.003}$ & $0.901\,{\scriptstyle\pm 0.003}$ & $0.900\,{\scriptstyle\pm 0.007}$ & -- & $0.051\,{\scriptstyle\pm 0.001}$ \\
H-Z & $0.897\,{\scriptstyle\pm 0.003}$ & $1.000\,{\scriptstyle\pm 0.000}$ & $0.588\,{\scriptstyle\pm 0.010}$ & -- & $0.561\,{\scriptstyle\pm 0.001}$ \\
H-GZ & $0.901\,{\scriptstyle\pm 0.003}$ & $0.901\,{\scriptstyle\pm 0.003}$ & $0.900\,{\scriptstyle\pm 0.007}$ & -- & $0.538\,{\scriptstyle\pm 0.001}$ \\
Oracle & $0.898\,{\scriptstyle\pm 0.003}$ & $0.972\,{\scriptstyle\pm 0.001}$ & $0.675\,{\scriptstyle\pm 0.008}$ & -- & $0.019\,{\scriptstyle\pm 0.000}$ \\
Coord & $0.824\,{\scriptstyle\pm 0.003}$ & $0.989\,{\scriptstyle\pm 0.001}$ & $0.329\,{\scriptstyle\pm 0.009}$ & -- & $0.039\,{\scriptstyle\pm 0.001}$ \\
Bonf & $0.930\,{\scriptstyle\pm 0.002}$ & $1.000\,{\scriptstyle\pm 0.000}$ & $0.722\,{\scriptstyle\pm 0.008}$ & -- & $0.121\,{\scriptstyle\pm 0.002}$ \\
H-weight & $0.897\,{\scriptstyle\pm 0.003}$ & $1.000\,{\scriptstyle\pm 0.000}$ & $0.588\,{\scriptstyle\pm 0.010}$ & -- & $0.070\,{\scriptstyle\pm 0.002}$ \\
CLR$_{10^{-2}}$ & $0.895\,{\scriptstyle\pm 0.003}$ & $1.000\,{\scriptstyle\pm 0.000}$ & $0.581\,{\scriptstyle\pm 0.010}$ & -- & $0.095\,{\scriptstyle\pm 0.002}$ \\
CLR$_{10^{-4}}$ & $0.896\,{\scriptstyle\pm 0.003}$ & $1.000\,{\scriptstyle\pm 0.000}$ & $0.586\,{\scriptstyle\pm 0.010}$ & -- & $0.130\,{\scriptstyle\pm 0.004}$ \\
\midrule
\multicolumn{6}{l}{$K=8$} \\
Max & $0.901\,{\scriptstyle\pm 0.003}$ & $1.000\,{\scriptstyle\pm 0.000}$ & $0.603\,{\scriptstyle\pm 0.011}$ & -- & -- \\
L2 & $0.902\,{\scriptstyle\pm 0.003}$ & $1.000\,{\scriptstyle\pm 0.000}$ & $0.608\,{\scriptstyle\pm 0.011}$ & -- & -- \\
H & $0.902\,{\scriptstyle\pm 0.003}$ & $1.000\,{\scriptstyle\pm 0.000}$ & $0.604\,{\scriptstyle\pm 0.011}$ & -- & -- \\
CLR & $0.902\,{\scriptstyle\pm 0.003}$ & $1.000\,{\scriptstyle\pm 0.000}$ & $0.606\,{\scriptstyle\pm 0.011}$ & -- & -- \\
CLR$_{10^{-3}}$ & $0.902\,{\scriptstyle\pm 0.003}$ & $1.000\,{\scriptstyle\pm 0.000}$ & $0.605\,{\scriptstyle\pm 0.011}$ & -- & -- \\
H-scale & $0.900\,{\scriptstyle\pm 0.003}$ & $0.895\,{\scriptstyle\pm 0.003}$ & $0.917\,{\scriptstyle\pm 0.004}$ & -- & -- \\
H-group & $0.903\,{\scriptstyle\pm 0.003}$ & $0.900\,{\scriptstyle\pm 0.004}$ & $0.910\,{\scriptstyle\pm 0.006}$ & -- & -- \\
H-Z & $0.902\,{\scriptstyle\pm 0.003}$ & $1.000\,{\scriptstyle\pm 0.000}$ & $0.604\,{\scriptstyle\pm 0.011}$ & -- & -- \\
H-GZ & $0.903\,{\scriptstyle\pm 0.003}$ & $0.900\,{\scriptstyle\pm 0.004}$ & $0.910\,{\scriptstyle\pm 0.006}$ & -- & -- \\
Oracle & $0.898\,{\scriptstyle\pm 0.002}$ & $0.981\,{\scriptstyle\pm 0.001}$ & $0.646\,{\scriptstyle\pm 0.007}$ & -- & -- \\
Coord & $0.697\,{\scriptstyle\pm 0.003}$ & $0.922\,{\scriptstyle\pm 0.003}$ & $0.020\,{\scriptstyle\pm 0.002}$ & -- & -- \\
Bonf & $0.912\,{\scriptstyle\pm 0.002}$ & $1.000\,{\scriptstyle\pm 0.000}$ & $0.644\,{\scriptstyle\pm 0.009}$ & -- & -- \\
H-weight & $0.902\,{\scriptstyle\pm 0.003}$ & $1.000\,{\scriptstyle\pm 0.000}$ & $0.604\,{\scriptstyle\pm 0.011}$ & -- & -- \\
CLR$_{10^{-2}}$ & $0.901\,{\scriptstyle\pm 0.003}$ & $1.000\,{\scriptstyle\pm 0.000}$ & $0.601\,{\scriptstyle\pm 0.011}$ & -- & -- \\
CLR$_{10^{-4}}$ & $0.901\,{\scriptstyle\pm 0.003}$ & $1.000\,{\scriptstyle\pm 0.000}$ & $0.604\,{\scriptstyle\pm 0.011}$ & -- & -- \\
\midrule
\multicolumn{6}{l}{$K=20$} \\
Max & $0.899\,{\scriptstyle\pm 0.003}$ & $1.000\,{\scriptstyle\pm 0.000}$ & $0.596\,{\scriptstyle\pm 0.011}$ & -- & -- \\
L2 & $0.901\,{\scriptstyle\pm 0.003}$ & $1.000\,{\scriptstyle\pm 0.000}$ & $0.604\,{\scriptstyle\pm 0.011}$ & -- & -- \\
H & $0.902\,{\scriptstyle\pm 0.003}$ & $1.000\,{\scriptstyle\pm 0.000}$ & $0.607\,{\scriptstyle\pm 0.011}$ & -- & -- \\
CLR & $0.902\,{\scriptstyle\pm 0.003}$ & $1.000\,{\scriptstyle\pm 0.000}$ & $0.609\,{\scriptstyle\pm 0.011}$ & -- & -- \\
CLR$_{10^{-3}}$ & $0.902\,{\scriptstyle\pm 0.003}$ & $1.000\,{\scriptstyle\pm 0.000}$ & $0.607\,{\scriptstyle\pm 0.011}$ & -- & -- \\
H-scale & $0.900\,{\scriptstyle\pm 0.003}$ & $0.895\,{\scriptstyle\pm 0.003}$ & $0.915\,{\scriptstyle\pm 0.004}$ & -- & -- \\
H-group & $0.903\,{\scriptstyle\pm 0.003}$ & $0.903\,{\scriptstyle\pm 0.003}$ & $0.904\,{\scriptstyle\pm 0.006}$ & -- & -- \\
H-Z & $0.902\,{\scriptstyle\pm 0.003}$ & $1.000\,{\scriptstyle\pm 0.000}$ & $0.607\,{\scriptstyle\pm 0.011}$ & -- & -- \\
H-GZ & $0.903\,{\scriptstyle\pm 0.003}$ & $0.903\,{\scriptstyle\pm 0.003}$ & $0.904\,{\scriptstyle\pm 0.006}$ & -- & -- \\
Oracle & $0.903\,{\scriptstyle\pm 0.003}$ & $0.879\,{\scriptstyle\pm 0.004}$ & $0.973\,{\scriptstyle\pm 0.001}$ & -- & -- \\
Coord & $0.443\,{\scriptstyle\pm 0.005}$ & $0.590\,{\scriptstyle\pm 0.006}$ & $0.000\,{\scriptstyle\pm 0.000}$ & -- & -- \\
Bonf & $0.924\,{\scriptstyle\pm 0.002}$ & $1.000\,{\scriptstyle\pm 0.000}$ & $0.693\,{\scriptstyle\pm 0.009}$ & -- & -- \\
H-weight & $0.902\,{\scriptstyle\pm 0.003}$ & $1.000\,{\scriptstyle\pm 0.000}$ & $0.607\,{\scriptstyle\pm 0.011}$ & -- & -- \\
CLR$_{10^{-2}}$ & $0.902\,{\scriptstyle\pm 0.003}$ & $1.000\,{\scriptstyle\pm 0.000}$ & $0.607\,{\scriptstyle\pm 0.012}$ & -- & -- \\
CLR$_{10^{-4}}$ & $0.901\,{\scriptstyle\pm 0.003}$ & $1.000\,{\scriptstyle\pm 0.000}$ & $0.602\,{\scriptstyle\pm 0.011}$ & -- & -- \\
\bottomrule
\end{tabular}
\end{center}
\clearpage
\subsection{Zeros}
\begin{center}\small\setlength{\tabcolsep}{4pt}\begin{tabular}{lccccc}\toprule
Method & Marginal & $G=0$ & $G=1$ & Boundary & $V_{1/2}$ \\\midrule
\multicolumn{6}{l}{$K=3$} \\
Max & $0.901\,{\scriptstyle\pm 0.003}$ & $0.901\,{\scriptstyle\pm 0.003}$ & $0.900\,{\scriptstyle\pm 0.004}$ & $0.410\,{\scriptstyle\pm 0.010}$ & $0.136\,{\scriptstyle\pm 0.004}$ \\
L2 & $0.901\,{\scriptstyle\pm 0.003}$ & $0.902\,{\scriptstyle\pm 0.003}$ & $0.900\,{\scriptstyle\pm 0.004}$ & $0.411\,{\scriptstyle\pm 0.010}$ & $0.140\,{\scriptstyle\pm 0.004}$ \\
H & $0.901\,{\scriptstyle\pm 0.003}$ & $0.902\,{\scriptstyle\pm 0.003}$ & $0.900\,{\scriptstyle\pm 0.004}$ & $0.343\,{\scriptstyle\pm 0.017}$ & $0.226\,{\scriptstyle\pm 0.011}$ \\
CLR & $1.000\,{\scriptstyle\pm 0.000}$ & $1.000\,{\scriptstyle\pm 0.000}$ & $1.000\,{\scriptstyle\pm 0.000}$ & $1.000\,{\scriptstyle\pm 0.000}$ & $1.000\,{\scriptstyle\pm 0.000}$ \\
CLR$_{10^{-3}}$ & $0.901\,{\scriptstyle\pm 0.003}$ & $0.902\,{\scriptstyle\pm 0.003}$ & $0.901\,{\scriptstyle\pm 0.004}$ & $0.343\,{\scriptstyle\pm 0.017}$ & $0.590\,{\scriptstyle\pm 0.007}$ \\
H-scale & $0.901\,{\scriptstyle\pm 0.003}$ & $0.902\,{\scriptstyle\pm 0.003}$ & $0.900\,{\scriptstyle\pm 0.004}$ & $0.343\,{\scriptstyle\pm 0.017}$ & $0.226\,{\scriptstyle\pm 0.010}$ \\
H-group & $0.903\,{\scriptstyle\pm 0.003}$ & $0.904\,{\scriptstyle\pm 0.003}$ & $0.900\,{\scriptstyle\pm 0.006}$ & $0.357\,{\scriptstyle\pm 0.017}$ & $0.235\,{\scriptstyle\pm 0.010}$ \\
H-Z & $0.902\,{\scriptstyle\pm 0.003}$ & $0.902\,{\scriptstyle\pm 0.003}$ & $0.901\,{\scriptstyle\pm 0.004}$ & $0.905\,{\scriptstyle\pm 0.007}$ & $0.405\,{\scriptstyle\pm 0.002}$ \\
H-GZ & $0.903\,{\scriptstyle\pm 0.003}$ & $0.902\,{\scriptstyle\pm 0.003}$ & $0.905\,{\scriptstyle\pm 0.005}$ & $0.913\,{\scriptstyle\pm 0.007}$ & $0.409\,{\scriptstyle\pm 0.002}$ \\
Oracle & $0.901\,{\scriptstyle\pm 0.003}$ & $0.902\,{\scriptstyle\pm 0.003}$ & $0.899\,{\scriptstyle\pm 0.004}$ & $0.459\,{\scriptstyle\pm 0.016}$ & $0.086\,{\scriptstyle\pm 0.002}$ \\
Coord & $0.833\,{\scriptstyle\pm 0.004}$ & $0.834\,{\scriptstyle\pm 0.005}$ & $0.831\,{\scriptstyle\pm 0.005}$ & $0.263\,{\scriptstyle\pm 0.007}$ & $0.075\,{\scriptstyle\pm 0.002}$ \\
Bonf & $0.946\,{\scriptstyle\pm 0.002}$ & $0.947\,{\scriptstyle\pm 0.002}$ & $0.945\,{\scriptstyle\pm 0.003}$ & $0.647\,{\scriptstyle\pm 0.010}$ & $0.315\,{\scriptstyle\pm 0.013}$ \\
H-weight & $0.901\,{\scriptstyle\pm 0.003}$ & $0.902\,{\scriptstyle\pm 0.003}$ & $0.900\,{\scriptstyle\pm 0.004}$ & $0.343\,{\scriptstyle\pm 0.017}$ & $0.226\,{\scriptstyle\pm 0.011}$ \\
CLR$_{10^{-2}}$ & $0.901\,{\scriptstyle\pm 0.003}$ & $0.902\,{\scriptstyle\pm 0.003}$ & $0.901\,{\scriptstyle\pm 0.004}$ & $0.343\,{\scriptstyle\pm 0.017}$ & $0.406\,{\scriptstyle\pm 0.013}$ \\
CLR$_{10^{-4}}$ & $0.901\,{\scriptstyle\pm 0.003}$ & $0.902\,{\scriptstyle\pm 0.003}$ & $0.901\,{\scriptstyle\pm 0.004}$ & $0.344\,{\scriptstyle\pm 0.017}$ & $0.627\,{\scriptstyle\pm 0.005}$ \\
\midrule
\multicolumn{6}{l}{$K=8$} \\
Max & $0.900\,{\scriptstyle\pm 0.003}$ & $0.900\,{\scriptstyle\pm 0.003}$ & $0.899\,{\scriptstyle\pm 0.004}$ & $0.639\,{\scriptstyle\pm 0.007}$ & -- \\
L2 & $0.900\,{\scriptstyle\pm 0.003}$ & $0.900\,{\scriptstyle\pm 0.003}$ & $0.900\,{\scriptstyle\pm 0.003}$ & $0.614\,{\scriptstyle\pm 0.006}$ & -- \\
H & $0.899\,{\scriptstyle\pm 0.003}$ & $0.900\,{\scriptstyle\pm 0.003}$ & $0.898\,{\scriptstyle\pm 0.004}$ & $0.409\,{\scriptstyle\pm 0.009}$ & -- \\
CLR & $1.000\,{\scriptstyle\pm 0.000}$ & $1.000\,{\scriptstyle\pm 0.000}$ & $1.000\,{\scriptstyle\pm 0.000}$ & $1.000\,{\scriptstyle\pm 0.000}$ & -- \\
CLR$_{10^{-3}}$ & $0.899\,{\scriptstyle\pm 0.003}$ & $0.899\,{\scriptstyle\pm 0.003}$ & $0.899\,{\scriptstyle\pm 0.003}$ & $0.329\,{\scriptstyle\pm 0.015}$ & -- \\
H-scale & $0.899\,{\scriptstyle\pm 0.003}$ & $0.900\,{\scriptstyle\pm 0.003}$ & $0.898\,{\scriptstyle\pm 0.004}$ & $0.410\,{\scriptstyle\pm 0.009}$ & -- \\
H-group & $0.901\,{\scriptstyle\pm 0.003}$ & $0.901\,{\scriptstyle\pm 0.003}$ & $0.901\,{\scriptstyle\pm 0.006}$ & $0.418\,{\scriptstyle\pm 0.009}$ & -- \\
H-Z & $0.902\,{\scriptstyle\pm 0.003}$ & $0.902\,{\scriptstyle\pm 0.003}$ & $0.902\,{\scriptstyle\pm 0.004}$ & $0.900\,{\scriptstyle\pm 0.008}$ & -- \\
H-GZ & $0.904\,{\scriptstyle\pm 0.003}$ & $0.904\,{\scriptstyle\pm 0.003}$ & $0.904\,{\scriptstyle\pm 0.005}$ & $0.904\,{\scriptstyle\pm 0.007}$ & -- \\
Oracle & $0.900\,{\scriptstyle\pm 0.003}$ & $0.900\,{\scriptstyle\pm 0.003}$ & $0.901\,{\scriptstyle\pm 0.004}$ & $0.685\,{\scriptstyle\pm 0.010}$ & -- \\
Coord & $0.450\,{\scriptstyle\pm 0.004}$ & $0.449\,{\scriptstyle\pm 0.004}$ & $0.452\,{\scriptstyle\pm 0.006}$ & $0.196\,{\scriptstyle\pm 0.005}$ & -- \\
Bonf & $0.916\,{\scriptstyle\pm 0.003}$ & $0.916\,{\scriptstyle\pm 0.003}$ & $0.916\,{\scriptstyle\pm 0.003}$ & $0.706\,{\scriptstyle\pm 0.009}$ & -- \\
H-weight & $0.899\,{\scriptstyle\pm 0.003}$ & $0.900\,{\scriptstyle\pm 0.003}$ & $0.898\,{\scriptstyle\pm 0.004}$ & $0.409\,{\scriptstyle\pm 0.009}$ & -- \\
CLR$_{10^{-2}}$ & $0.900\,{\scriptstyle\pm 0.003}$ & $0.900\,{\scriptstyle\pm 0.003}$ & $0.899\,{\scriptstyle\pm 0.003}$ & $0.384\,{\scriptstyle\pm 0.010}$ & -- \\
CLR$_{10^{-4}}$ & $0.899\,{\scriptstyle\pm 0.003}$ & $0.899\,{\scriptstyle\pm 0.003}$ & $0.899\,{\scriptstyle\pm 0.003}$ & $0.324\,{\scriptstyle\pm 0.016}$ & -- \\
\midrule
\multicolumn{6}{l}{$K=20$} \\
Max & $0.898\,{\scriptstyle\pm 0.003}$ & $0.897\,{\scriptstyle\pm 0.003}$ & $0.899\,{\scriptstyle\pm 0.004}$ & $0.841\,{\scriptstyle\pm 0.006}$ & -- \\
L2 & $0.900\,{\scriptstyle\pm 0.003}$ & $0.899\,{\scriptstyle\pm 0.003}$ & $0.901\,{\scriptstyle\pm 0.004}$ & $0.791\,{\scriptstyle\pm 0.006}$ & -- \\
H & $0.904\,{\scriptstyle\pm 0.003}$ & $0.903\,{\scriptstyle\pm 0.003}$ & $0.905\,{\scriptstyle\pm 0.003}$ & $0.621\,{\scriptstyle\pm 0.007}$ & -- \\
CLR & $1.000\,{\scriptstyle\pm 0.000}$ & $1.000\,{\scriptstyle\pm 0.000}$ & $1.000\,{\scriptstyle\pm 0.000}$ & $1.000\,{\scriptstyle\pm 0.000}$ & -- \\
CLR$_{10^{-3}}$ & $0.903\,{\scriptstyle\pm 0.003}$ & $0.903\,{\scriptstyle\pm 0.003}$ & $0.903\,{\scriptstyle\pm 0.004}$ & $0.529\,{\scriptstyle\pm 0.009}$ & -- \\
H-scale & $0.904\,{\scriptstyle\pm 0.003}$ & $0.904\,{\scriptstyle\pm 0.003}$ & $0.903\,{\scriptstyle\pm 0.004}$ & $0.621\,{\scriptstyle\pm 0.008}$ & -- \\
H-group & $0.905\,{\scriptstyle\pm 0.003}$ & $0.904\,{\scriptstyle\pm 0.003}$ & $0.909\,{\scriptstyle\pm 0.006}$ & $0.625\,{\scriptstyle\pm 0.008}$ & -- \\
H-Z & $0.904\,{\scriptstyle\pm 0.003}$ & $0.903\,{\scriptstyle\pm 0.003}$ & $0.905\,{\scriptstyle\pm 0.004}$ & $0.907\,{\scriptstyle\pm 0.007}$ & -- \\
H-GZ & $0.905\,{\scriptstyle\pm 0.003}$ & $0.903\,{\scriptstyle\pm 0.004}$ & $0.912\,{\scriptstyle\pm 0.006}$ & $0.913\,{\scriptstyle\pm 0.008}$ & -- \\
Oracle & $0.903\,{\scriptstyle\pm 0.003}$ & $0.903\,{\scriptstyle\pm 0.003}$ & $0.904\,{\scriptstyle\pm 0.004}$ & $0.813\,{\scriptstyle\pm 0.006}$ & -- \\
Coord & $0.083\,{\scriptstyle\pm 0.002}$ & $0.083\,{\scriptstyle\pm 0.002}$ & $0.083\,{\scriptstyle\pm 0.003}$ & $0.050\,{\scriptstyle\pm 0.002}$ & -- \\
Bonf & $0.921\,{\scriptstyle\pm 0.002}$ & $0.921\,{\scriptstyle\pm 0.002}$ & $0.920\,{\scriptstyle\pm 0.003}$ & $0.894\,{\scriptstyle\pm 0.004}$ & -- \\
H-weight & $0.904\,{\scriptstyle\pm 0.003}$ & $0.903\,{\scriptstyle\pm 0.003}$ & $0.905\,{\scriptstyle\pm 0.003}$ & $0.621\,{\scriptstyle\pm 0.007}$ & -- \\
CLR$_{10^{-2}}$ & $0.903\,{\scriptstyle\pm 0.003}$ & $0.903\,{\scriptstyle\pm 0.003}$ & $0.904\,{\scriptstyle\pm 0.004}$ & $0.631\,{\scriptstyle\pm 0.007}$ & -- \\
CLR$_{10^{-4}}$ & $0.901\,{\scriptstyle\pm 0.003}$ & $0.900\,{\scriptstyle\pm 0.003}$ & $0.902\,{\scriptstyle\pm 0.004}$ & $0.521\,{\scriptstyle\pm 0.011}$ & -- \\
\bottomrule
\end{tabular}
\end{center}
\clearpage
\subsection{Mixed}
\begin{center}\small\setlength{\tabcolsep}{4pt}\begin{tabular}{lccccc}\toprule
Method & Marginal & $G=0$ & $G=1$ & Boundary & $V_{1/2}$ \\\midrule
\multicolumn{6}{l}{$K=3$} \\
Max & $0.898\,{\scriptstyle\pm 0.003}$ & $0.989\,{\scriptstyle\pm 0.001}$ & $0.624\,{\scriptstyle\pm 0.010}$ & $0.331\,{\scriptstyle\pm 0.010}$ & $0.152\,{\scriptstyle\pm 0.004}$ \\
L2 & $0.898\,{\scriptstyle\pm 0.003}$ & $0.989\,{\scriptstyle\pm 0.001}$ & $0.624\,{\scriptstyle\pm 0.010}$ & $0.329\,{\scriptstyle\pm 0.010}$ & $0.152\,{\scriptstyle\pm 0.004}$ \\
H & $0.898\,{\scriptstyle\pm 0.003}$ & $0.985\,{\scriptstyle\pm 0.001}$ & $0.637\,{\scriptstyle\pm 0.011}$ & $0.187\,{\scriptstyle\pm 0.011}$ & $0.153\,{\scriptstyle\pm 0.006}$ \\
CLR & $0.927\,{\scriptstyle\pm 0.011}$ & $0.987\,{\scriptstyle\pm 0.002}$ & $0.749\,{\scriptstyle\pm 0.037}$ & $0.340\,{\scriptstyle\pm 0.094}$ & $0.625\,{\scriptstyle\pm 0.055}$ \\
CLR$_{10^{-3}}$ & $0.897\,{\scriptstyle\pm 0.003}$ & $0.981\,{\scriptstyle\pm 0.001}$ & $0.645\,{\scriptstyle\pm 0.012}$ & $0.054\,{\scriptstyle\pm 0.009}$ & $0.389\,{\scriptstyle\pm 0.012}$ \\
H-scale & $0.897\,{\scriptstyle\pm 0.003}$ & $0.942\,{\scriptstyle\pm 0.004}$ & $0.761\,{\scriptstyle\pm 0.010}$ & $0.283\,{\scriptstyle\pm 0.016}$ & $0.074\,{\scriptstyle\pm 0.003}$ \\
H-group & $0.899\,{\scriptstyle\pm 0.003}$ & $0.899\,{\scriptstyle\pm 0.003}$ & $0.901\,{\scriptstyle\pm 0.006}$ & $0.585\,{\scriptstyle\pm 0.015}$ & $0.178\,{\scriptstyle\pm 0.010}$ \\
H-Z & $0.900\,{\scriptstyle\pm 0.003}$ & $1.000\,{\scriptstyle\pm 0.000}$ & $0.602\,{\scriptstyle\pm 0.012}$ & $0.910\,{\scriptstyle\pm 0.008}$ & $0.460\,{\scriptstyle\pm 0.004}$ \\
H-GZ & $0.903\,{\scriptstyle\pm 0.003}$ & $0.901\,{\scriptstyle\pm 0.004}$ & $0.911\,{\scriptstyle\pm 0.005}$ & $0.921\,{\scriptstyle\pm 0.007}$ & $0.490\,{\scriptstyle\pm 0.009}$ \\
Oracle & $0.900\,{\scriptstyle\pm 0.003}$ & $0.972\,{\scriptstyle\pm 0.001}$ & $0.683\,{\scriptstyle\pm 0.010}$ & $0.279\,{\scriptstyle\pm 0.015}$ & $0.038\,{\scriptstyle\pm 0.001}$ \\
Coord & $0.828\,{\scriptstyle\pm 0.003}$ & $0.983\,{\scriptstyle\pm 0.001}$ & $0.363\,{\scriptstyle\pm 0.011}$ & $0.154\,{\scriptstyle\pm 0.006}$ & $0.063\,{\scriptstyle\pm 0.002}$ \\
Bonf & $0.942\,{\scriptstyle\pm 0.002}$ & $0.992\,{\scriptstyle\pm 0.001}$ & $0.790\,{\scriptstyle\pm 0.009}$ & $0.566\,{\scriptstyle\pm 0.015}$ & $0.245\,{\scriptstyle\pm 0.009}$ \\
H-weight & $0.898\,{\scriptstyle\pm 0.003}$ & $0.985\,{\scriptstyle\pm 0.001}$ & $0.637\,{\scriptstyle\pm 0.011}$ & $0.187\,{\scriptstyle\pm 0.011}$ & $0.153\,{\scriptstyle\pm 0.006}$ \\
CLR$_{10^{-2}}$ & $0.896\,{\scriptstyle\pm 0.003}$ & $0.983\,{\scriptstyle\pm 0.001}$ & $0.638\,{\scriptstyle\pm 0.011}$ & $0.134\,{\scriptstyle\pm 0.009}$ & $0.218\,{\scriptstyle\pm 0.008}$ \\
CLR$_{10^{-4}}$ & $0.896\,{\scriptstyle\pm 0.003}$ & $0.981\,{\scriptstyle\pm 0.001}$ & $0.644\,{\scriptstyle\pm 0.011}$ & $0.028\,{\scriptstyle\pm 0.008}$ & $0.454\,{\scriptstyle\pm 0.015}$ \\
\midrule
\multicolumn{6}{l}{$K=8$} \\
Max & $0.900\,{\scriptstyle\pm 0.003}$ & $0.995\,{\scriptstyle\pm 0.000}$ & $0.611\,{\scriptstyle\pm 0.011}$ & $0.515\,{\scriptstyle\pm 0.011}$ & -- \\
L2 & $0.899\,{\scriptstyle\pm 0.003}$ & $0.996\,{\scriptstyle\pm 0.000}$ & $0.604\,{\scriptstyle\pm 0.011}$ & $0.509\,{\scriptstyle\pm 0.012}$ & -- \\
H & $0.898\,{\scriptstyle\pm 0.003}$ & $0.995\,{\scriptstyle\pm 0.000}$ & $0.607\,{\scriptstyle\pm 0.012}$ & $0.413\,{\scriptstyle\pm 0.012}$ & -- \\
CLR & $0.931\,{\scriptstyle\pm 0.011}$ & $0.987\,{\scriptstyle\pm 0.002}$ & $0.762\,{\scriptstyle\pm 0.037}$ & $0.370\,{\scriptstyle\pm 0.096}$ & -- \\
CLR$_{10^{-3}}$ & $0.899\,{\scriptstyle\pm 0.003}$ & $0.991\,{\scriptstyle\pm 0.001}$ & $0.620\,{\scriptstyle\pm 0.012}$ & $0.300\,{\scriptstyle\pm 0.012}$ & -- \\
H-scale & $0.899\,{\scriptstyle\pm 0.003}$ & $0.917\,{\scriptstyle\pm 0.003}$ & $0.842\,{\scriptstyle\pm 0.006}$ & $0.531\,{\scriptstyle\pm 0.010}$ & -- \\
H-group & $0.901\,{\scriptstyle\pm 0.003}$ & $0.900\,{\scriptstyle\pm 0.004}$ & $0.902\,{\scriptstyle\pm 0.006}$ & $0.624\,{\scriptstyle\pm 0.012}$ & -- \\
H-Z & $0.901\,{\scriptstyle\pm 0.003}$ & $1.000\,{\scriptstyle\pm 0.000}$ & $0.603\,{\scriptstyle\pm 0.011}$ & $0.908\,{\scriptstyle\pm 0.011}$ & -- \\
H-GZ & $0.904\,{\scriptstyle\pm 0.003}$ & $0.903\,{\scriptstyle\pm 0.004}$ & $0.906\,{\scriptstyle\pm 0.006}$ & $0.919\,{\scriptstyle\pm 0.010}$ & -- \\
Oracle & $0.901\,{\scriptstyle\pm 0.003}$ & $0.984\,{\scriptstyle\pm 0.001}$ & $0.649\,{\scriptstyle\pm 0.009}$ & $0.578\,{\scriptstyle\pm 0.012}$ & -- \\
Coord & $0.694\,{\scriptstyle\pm 0.003}$ & $0.917\,{\scriptstyle\pm 0.003}$ & $0.022\,{\scriptstyle\pm 0.002}$ & $0.055\,{\scriptstyle\pm 0.004}$ & -- \\
Bonf & $0.911\,{\scriptstyle\pm 0.002}$ & $0.999\,{\scriptstyle\pm 0.000}$ & $0.647\,{\scriptstyle\pm 0.009}$ & $0.605\,{\scriptstyle\pm 0.011}$ & -- \\
H-weight & $0.898\,{\scriptstyle\pm 0.003}$ & $0.995\,{\scriptstyle\pm 0.000}$ & $0.607\,{\scriptstyle\pm 0.012}$ & $0.413\,{\scriptstyle\pm 0.012}$ & -- \\
CLR$_{10^{-2}}$ & $0.899\,{\scriptstyle\pm 0.003}$ & $0.992\,{\scriptstyle\pm 0.001}$ & $0.616\,{\scriptstyle\pm 0.012}$ & $0.392\,{\scriptstyle\pm 0.012}$ & -- \\
CLR$_{10^{-4}}$ & $0.898\,{\scriptstyle\pm 0.003}$ & $0.990\,{\scriptstyle\pm 0.001}$ & $0.620\,{\scriptstyle\pm 0.011}$ & $0.244\,{\scriptstyle\pm 0.013}$ & -- \\
\midrule
\multicolumn{6}{l}{$K=20$} \\
Max & $0.899\,{\scriptstyle\pm 0.003}$ & $1.000\,{\scriptstyle\pm 0.000}$ & $0.594\,{\scriptstyle\pm 0.012}$ & $0.633\,{\scriptstyle\pm 0.011}$ & -- \\
L2 & $0.899\,{\scriptstyle\pm 0.003}$ & $1.000\,{\scriptstyle\pm 0.000}$ & $0.598\,{\scriptstyle\pm 0.012}$ & $0.619\,{\scriptstyle\pm 0.013}$ & -- \\
H & $0.900\,{\scriptstyle\pm 0.003}$ & $1.000\,{\scriptstyle\pm 0.000}$ & $0.601\,{\scriptstyle\pm 0.013}$ & $0.574\,{\scriptstyle\pm 0.014}$ & -- \\
CLR & $0.936\,{\scriptstyle\pm 0.011}$ & $0.988\,{\scriptstyle\pm 0.002}$ & $0.778\,{\scriptstyle\pm 0.037}$ & $0.400\,{\scriptstyle\pm 0.098}$ & -- \\
CLR$_{10^{-3}}$ & $0.901\,{\scriptstyle\pm 0.003}$ & $1.000\,{\scriptstyle\pm 0.000}$ & $0.604\,{\scriptstyle\pm 0.013}$ & $0.555\,{\scriptstyle\pm 0.013}$ & -- \\
H-scale & $0.902\,{\scriptstyle\pm 0.003}$ & $0.898\,{\scriptstyle\pm 0.003}$ & $0.912\,{\scriptstyle\pm 0.004}$ & $0.764\,{\scriptstyle\pm 0.007}$ & -- \\
H-group & $0.904\,{\scriptstyle\pm 0.003}$ & $0.903\,{\scriptstyle\pm 0.003}$ & $0.905\,{\scriptstyle\pm 0.005}$ & $0.757\,{\scriptstyle\pm 0.007}$ & -- \\
H-Z & $0.902\,{\scriptstyle\pm 0.003}$ & $1.000\,{\scriptstyle\pm 0.000}$ & $0.608\,{\scriptstyle\pm 0.012}$ & $0.914\,{\scriptstyle\pm 0.009}$ & -- \\
H-GZ & $0.906\,{\scriptstyle\pm 0.003}$ & $0.905\,{\scriptstyle\pm 0.004}$ & $0.909\,{\scriptstyle\pm 0.005}$ & $0.923\,{\scriptstyle\pm 0.008}$ & -- \\
Oracle & $0.901\,{\scriptstyle\pm 0.003}$ & $0.878\,{\scriptstyle\pm 0.004}$ & $0.969\,{\scriptstyle\pm 0.002}$ & $0.873\,{\scriptstyle\pm 0.004}$ & -- \\
Coord & $0.438\,{\scriptstyle\pm 0.006}$ & $0.584\,{\scriptstyle\pm 0.008}$ & $0.000\,{\scriptstyle\pm 0.000}$ & $0.046\,{\scriptstyle\pm 0.003}$ & -- \\
Bonf & $0.922\,{\scriptstyle\pm 0.003}$ & $1.000\,{\scriptstyle\pm 0.000}$ & $0.690\,{\scriptstyle\pm 0.010}$ & $0.721\,{\scriptstyle\pm 0.011}$ & -- \\
H-weight & $0.900\,{\scriptstyle\pm 0.003}$ & $1.000\,{\scriptstyle\pm 0.000}$ & $0.601\,{\scriptstyle\pm 0.013}$ & $0.574\,{\scriptstyle\pm 0.014}$ & -- \\
CLR$_{10^{-2}}$ & $0.902\,{\scriptstyle\pm 0.003}$ & $1.000\,{\scriptstyle\pm 0.000}$ & $0.607\,{\scriptstyle\pm 0.012}$ & $0.604\,{\scriptstyle\pm 0.013}$ & -- \\
CLR$_{10^{-4}}$ & $0.901\,{\scriptstyle\pm 0.003}$ & $1.000\,{\scriptstyle\pm 0.000}$ & $0.605\,{\scriptstyle\pm 0.013}$ & $0.514\,{\scriptstyle\pm 0.013}$ & -- \\
\bottomrule
\end{tabular}
\end{center}
\clearpage
\subsection{Shift}
\begin{center}\small\setlength{\tabcolsep}{4pt}\begin{tabular}{lccccc}\toprule
Method & Marginal & $G=0$ & $G=1$ & Boundary & $V_{1/2}$ \\\midrule
\multicolumn{6}{l}{$K=3$} \\
Max & $0.756\,{\scriptstyle\pm 0.006}$ & $0.990\,{\scriptstyle\pm 0.001}$ & $0.630\,{\scriptstyle\pm 0.010}$ & $0.317\,{\scriptstyle\pm 0.010}$ & $0.154\,{\scriptstyle\pm 0.004}$ \\
L2 & $0.757\,{\scriptstyle\pm 0.007}$ & $0.989\,{\scriptstyle\pm 0.001}$ & $0.631\,{\scriptstyle\pm 0.010}$ & $0.318\,{\scriptstyle\pm 0.011}$ & $0.156\,{\scriptstyle\pm 0.005}$ \\
H & $0.764\,{\scriptstyle\pm 0.007}$ & $0.985\,{\scriptstyle\pm 0.001}$ & $0.644\,{\scriptstyle\pm 0.010}$ & $0.180\,{\scriptstyle\pm 0.010}$ & $0.157\,{\scriptstyle\pm 0.006}$ \\
CLR & $0.843\,{\scriptstyle\pm 0.025}$ & $0.988\,{\scriptstyle\pm 0.002}$ & $0.764\,{\scriptstyle\pm 0.038}$ & $0.380\,{\scriptstyle\pm 0.097}$ & $0.649\,{\scriptstyle\pm 0.056}$ \\
CLR$_{10^{-3}}$ & $0.768\,{\scriptstyle\pm 0.008}$ & $0.981\,{\scriptstyle\pm 0.001}$ & $0.653\,{\scriptstyle\pm 0.012}$ & $0.067\,{\scriptstyle\pm 0.012}$ & $0.399\,{\scriptstyle\pm 0.014}$ \\
H-scale & $0.824\,{\scriptstyle\pm 0.006}$ & $0.946\,{\scriptstyle\pm 0.004}$ & $0.758\,{\scriptstyle\pm 0.009}$ & $0.318\,{\scriptstyle\pm 0.017}$ & $0.171\,{\scriptstyle\pm 0.006}$ \\
H-group & $0.903\,{\scriptstyle\pm 0.004}$ & $0.902\,{\scriptstyle\pm 0.004}$ & $0.903\,{\scriptstyle\pm 0.006}$ & $0.679\,{\scriptstyle\pm 0.016}$ & $0.452\,{\scriptstyle\pm 0.027}$ \\
H-Z & $0.750\,{\scriptstyle\pm 0.007}$ & $1.000\,{\scriptstyle\pm 0.000}$ & $0.615\,{\scriptstyle\pm 0.010}$ & $0.900\,{\scriptstyle\pm 0.010}$ & $0.462\,{\scriptstyle\pm 0.004}$ \\
H-GZ & $0.906\,{\scriptstyle\pm 0.004}$ & $0.905\,{\scriptstyle\pm 0.004}$ & $0.906\,{\scriptstyle\pm 0.005}$ & $0.916\,{\scriptstyle\pm 0.008}$ & $0.520\,{\scriptstyle\pm 0.005}$ \\
Oracle & $0.785\,{\scriptstyle\pm 0.007}$ & $0.972\,{\scriptstyle\pm 0.001}$ & $0.684\,{\scriptstyle\pm 0.010}$ & $0.324\,{\scriptstyle\pm 0.019}$ & $0.069\,{\scriptstyle\pm 0.003}$ \\
Coord & $0.588\,{\scriptstyle\pm 0.007}$ & $0.984\,{\scriptstyle\pm 0.001}$ & $0.374\,{\scriptstyle\pm 0.010}$ & $0.141\,{\scriptstyle\pm 0.006}$ & $0.065\,{\scriptstyle\pm 0.002}$ \\
Bonf & $0.861\,{\scriptstyle\pm 0.006}$ & $0.992\,{\scriptstyle\pm 0.001}$ & $0.791\,{\scriptstyle\pm 0.009}$ & $0.557\,{\scriptstyle\pm 0.015}$ & $0.245\,{\scriptstyle\pm 0.009}$ \\
H-weight & $0.900\,{\scriptstyle\pm 0.005}$ & $0.993\,{\scriptstyle\pm 0.001}$ & $0.851\,{\scriptstyle\pm 0.007}$ & $0.553\,{\scriptstyle\pm 0.019}$ & $0.423\,{\scriptstyle\pm 0.032}$ \\
CLR$_{10^{-2}}$ & $0.768\,{\scriptstyle\pm 0.007}$ & $0.983\,{\scriptstyle\pm 0.001}$ & $0.651\,{\scriptstyle\pm 0.010}$ & $0.143\,{\scriptstyle\pm 0.010}$ & $0.230\,{\scriptstyle\pm 0.009}$ \\
CLR$_{10^{-4}}$ & $0.769\,{\scriptstyle\pm 0.008}$ & $0.981\,{\scriptstyle\pm 0.001}$ & $0.654\,{\scriptstyle\pm 0.012}$ & $0.048\,{\scriptstyle\pm 0.012}$ & $0.464\,{\scriptstyle\pm 0.016}$ \\
\midrule
\multicolumn{6}{l}{$K=8$} \\
Max & $0.749\,{\scriptstyle\pm 0.006}$ & $0.995\,{\scriptstyle\pm 0.000}$ & $0.615\,{\scriptstyle\pm 0.010}$ & $0.481\,{\scriptstyle\pm 0.009}$ & -- \\
L2 & $0.747\,{\scriptstyle\pm 0.006}$ & $0.997\,{\scriptstyle\pm 0.000}$ & $0.612\,{\scriptstyle\pm 0.009}$ & $0.473\,{\scriptstyle\pm 0.009}$ & -- \\
H & $0.751\,{\scriptstyle\pm 0.007}$ & $0.995\,{\scriptstyle\pm 0.000}$ & $0.619\,{\scriptstyle\pm 0.011}$ & $0.379\,{\scriptstyle\pm 0.011}$ & -- \\
CLR & $0.840\,{\scriptstyle\pm 0.024}$ & $0.987\,{\scriptstyle\pm 0.002}$ & $0.760\,{\scriptstyle\pm 0.036}$ & $0.350\,{\scriptstyle\pm 0.095}$ & -- \\
CLR$_{10^{-3}}$ & $0.756\,{\scriptstyle\pm 0.007}$ & $0.991\,{\scriptstyle\pm 0.001}$ & $0.629\,{\scriptstyle\pm 0.010}$ & $0.271\,{\scriptstyle\pm 0.010}$ & -- \\
H-scale & $0.870\,{\scriptstyle\pm 0.003}$ & $0.921\,{\scriptstyle\pm 0.003}$ & $0.843\,{\scriptstyle\pm 0.005}$ & $0.608\,{\scriptstyle\pm 0.008}$ & -- \\
H-group & $0.904\,{\scriptstyle\pm 0.003}$ & $0.901\,{\scriptstyle\pm 0.003}$ & $0.905\,{\scriptstyle\pm 0.005}$ & $0.718\,{\scriptstyle\pm 0.010}$ & -- \\
H-Z & $0.741\,{\scriptstyle\pm 0.007}$ & $1.000\,{\scriptstyle\pm 0.000}$ & $0.601\,{\scriptstyle\pm 0.011}$ & $0.898\,{\scriptstyle\pm 0.008}$ & -- \\
H-GZ & $0.906\,{\scriptstyle\pm 0.003}$ & $0.904\,{\scriptstyle\pm 0.004}$ & $0.907\,{\scriptstyle\pm 0.005}$ & $0.912\,{\scriptstyle\pm 0.008}$ & -- \\
Oracle & $0.761\,{\scriptstyle\pm 0.006}$ & $0.983\,{\scriptstyle\pm 0.002}$ & $0.641\,{\scriptstyle\pm 0.008}$ & $0.559\,{\scriptstyle\pm 0.010}$ & -- \\
Coord & $0.336\,{\scriptstyle\pm 0.002}$ & $0.919\,{\scriptstyle\pm 0.003}$ & $0.021\,{\scriptstyle\pm 0.002}$ & $0.019\,{\scriptstyle\pm 0.002}$ & -- \\
Bonf & $0.774\,{\scriptstyle\pm 0.005}$ & $0.999\,{\scriptstyle\pm 0.000}$ & $0.653\,{\scriptstyle\pm 0.008}$ & $0.568\,{\scriptstyle\pm 0.010}$ & -- \\
H-weight & $0.903\,{\scriptstyle\pm 0.004}$ & $0.997\,{\scriptstyle\pm 0.000}$ & $0.851\,{\scriptstyle\pm 0.006}$ & $0.651\,{\scriptstyle\pm 0.010}$ & -- \\
CLR$_{10^{-2}}$ & $0.750\,{\scriptstyle\pm 0.007}$ & $0.992\,{\scriptstyle\pm 0.001}$ & $0.618\,{\scriptstyle\pm 0.010}$ & $0.369\,{\scriptstyle\pm 0.010}$ & -- \\
CLR$_{10^{-4}}$ & $0.759\,{\scriptstyle\pm 0.007}$ & $0.990\,{\scriptstyle\pm 0.001}$ & $0.634\,{\scriptstyle\pm 0.010}$ & $0.223\,{\scriptstyle\pm 0.011}$ & -- \\
\midrule
\multicolumn{6}{l}{$K=20$} \\
Max & $0.736\,{\scriptstyle\pm 0.007}$ & $1.000\,{\scriptstyle\pm 0.000}$ & $0.595\,{\scriptstyle\pm 0.011}$ & $0.576\,{\scriptstyle\pm 0.012}$ & -- \\
L2 & $0.736\,{\scriptstyle\pm 0.006}$ & $1.000\,{\scriptstyle\pm 0.000}$ & $0.594\,{\scriptstyle\pm 0.010}$ & $0.560\,{\scriptstyle\pm 0.011}$ & -- \\
H & $0.737\,{\scriptstyle\pm 0.007}$ & $1.000\,{\scriptstyle\pm 0.000}$ & $0.596\,{\scriptstyle\pm 0.011}$ & $0.508\,{\scriptstyle\pm 0.012}$ & -- \\
CLR & $0.863\,{\scriptstyle\pm 0.025}$ & $0.990\,{\scriptstyle\pm 0.002}$ & $0.796\,{\scriptstyle\pm 0.038}$ & $0.460\,{\scriptstyle\pm 0.099}$ & -- \\
CLR$_{10^{-3}}$ & $0.738\,{\scriptstyle\pm 0.007}$ & $1.000\,{\scriptstyle\pm 0.000}$ & $0.597\,{\scriptstyle\pm 0.011}$ & $0.481\,{\scriptstyle\pm 0.012}$ & -- \\
H-scale & $0.900\,{\scriptstyle\pm 0.003}$ & $0.889\,{\scriptstyle\pm 0.004}$ & $0.906\,{\scriptstyle\pm 0.003}$ & $0.835\,{\scriptstyle\pm 0.005}$ & -- \\
H-group & $0.897\,{\scriptstyle\pm 0.003}$ & $0.894\,{\scriptstyle\pm 0.004}$ & $0.898\,{\scriptstyle\pm 0.005}$ & $0.825\,{\scriptstyle\pm 0.007}$ & -- \\
H-Z & $0.738\,{\scriptstyle\pm 0.007}$ & $1.000\,{\scriptstyle\pm 0.000}$ & $0.598\,{\scriptstyle\pm 0.010}$ & $0.890\,{\scriptstyle\pm 0.008}$ & -- \\
H-GZ & $0.902\,{\scriptstyle\pm 0.003}$ & $0.896\,{\scriptstyle\pm 0.004}$ & $0.905\,{\scriptstyle\pm 0.005}$ & $0.910\,{\scriptstyle\pm 0.007}$ & -- \\
Oracle & $0.935\,{\scriptstyle\pm 0.002}$ & $0.873\,{\scriptstyle\pm 0.005}$ & $0.969\,{\scriptstyle\pm 0.001}$ & $0.944\,{\scriptstyle\pm 0.002}$ & -- \\
Coord & $0.203\,{\scriptstyle\pm 0.003}$ & $0.581\,{\scriptstyle\pm 0.007}$ & $0.000\,{\scriptstyle\pm 0.000}$ & $0.009\,{\scriptstyle\pm 0.001}$ & -- \\
Bonf & $0.802\,{\scriptstyle\pm 0.006}$ & $1.000\,{\scriptstyle\pm 0.000}$ & $0.696\,{\scriptstyle\pm 0.010}$ & $0.682\,{\scriptstyle\pm 0.011}$ & -- \\
H-weight & $0.900\,{\scriptstyle\pm 0.004}$ & $1.000\,{\scriptstyle\pm 0.000}$ & $0.847\,{\scriptstyle\pm 0.006}$ & $0.785\,{\scriptstyle\pm 0.009}$ & -- \\
CLR$_{10^{-2}}$ & $0.740\,{\scriptstyle\pm 0.007}$ & $1.000\,{\scriptstyle\pm 0.000}$ & $0.601\,{\scriptstyle\pm 0.010}$ & $0.533\,{\scriptstyle\pm 0.011}$ & -- \\
CLR$_{10^{-4}}$ & $0.742\,{\scriptstyle\pm 0.008}$ & $1.000\,{\scriptstyle\pm 0.000}$ & $0.604\,{\scriptstyle\pm 0.011}$ & $0.438\,{\scriptstyle\pm 0.012}$ & -- \\
\bottomrule
\end{tabular}
\end{center}

\clearpage
\section*{NeurIPS Paper Checklist}




\begin{enumerate}

\item {\bf Claims}
    \item[] Question: Do the main claims made in the abstract and introduction accurately reflect the paper's contributions and scope?
    \item[] Answer: \answerYes{} % Replace by \answerYes{}, \answerNo{}, or \answerNA{}.
    \item[] Justification: Sections 1 and 6 state the synthetic scope, distinct reliability objectives, and exact score-extension convention. Numerical claims are generated from executed runs.
    \item[] Guidelines:
    \begin{itemize}
        \item The answer \answerNA{} means that the abstract and introduction do not include the claims made in the paper.
        \item The abstract and/or introduction should clearly state the claims made, including the contributions made in the paper and important assumptions and limitations. A \answerNo{} or \answerNA{} answer to this question will not be perceived well by the reviewers. 
        \item The claims made should match theoretical and experimental results, and reflect how much the results can be expected to generalize to other settings. 
        \item It is fine to include aspirational goals as motivation as long as it is clear that these goals are not attained by the paper. 
    \end{itemize}

\item {\bf Limitations}
    \item[] Question: Does the paper discuss the limitations of the work performed by the authors?
    \item[] Answer: \answerYes{} % Replace by \answerYes{}, \answerNo{}, or \answerNA{}.
    \item[] Justification: Section 7 discusses synthetic-only evidence, the fixed predictor, limited support mechanisms, high-dimensional size, and exploratory additions.
    \item[] Guidelines:
    \begin{itemize}
        \item The answer \answerNA{} means that the paper has no limitation while the answer \answerNo{} means that the paper has limitations, but those are not discussed in the paper. 
        \item The authors are encouraged to create a separate ``Limitations'' section in their paper.
        \item The paper should point out any strong assumptions and how robust the results are to violations of these assumptions (e.g., independence assumptions, noiseless settings, model well-specification, asymptotic approximations only holding locally). The authors should reflect on how these assumptions might be violated in practice and what the implications would be.
        \item The authors should reflect on the scope of the claims made, e.g., if the approach was only tested on a few datasets or with a few runs. In general, empirical results often depend on implicit assumptions, which should be articulated.
        \item The authors should reflect on the factors that influence the performance of the approach. For example, a facial recognition algorithm may perform poorly when image resolution is low or images are taken in low lighting. Or a speech-to-text system might not be used reliably to provide closed captions for online lectures because it fails to handle technical jargon.
        \item The authors should discuss the computational efficiency of the proposed algorithms and how they scale with dataset size.
        \item If applicable, the authors should discuss possible limitations of their approach to address problems of privacy and fairness.
        \item While the authors might fear that complete honesty about limitations might be used by reviewers as grounds for rejection, a worse outcome might be that reviewers discover limitations that aren't acknowledged in the paper. The authors should use their best judgment and recognize that individual actions in favor of transparency play an important role in developing norms that preserve the integrity of the community. Reviewers will be specifically instructed to not penalize honesty concerning limitations.
    \end{itemize}

\item {\bf Theory assumptions and proofs}
    \item[] Question: For each theoretical result, does the paper provide the full set of assumptions and a complete (and correct) proof?
    \item[] Answer: \answerYes{} % Replace by \answerYes{}, \answerNo{}, or \answerNA{}.
    \item[] Justification: Section 4 and Appendix A state the independence, tie, and support assumptions and provide complete proofs. Appendix B specifies the oracle efficiency target.
    \item[] Guidelines:
    \begin{itemize}
        \item The answer \answerNA{} means that the paper does not include theoretical results. 
        \item All the theorems, formulas, and proofs in the paper should be numbered and cross-referenced.
        \item All assumptions should be clearly stated or referenced in the statement of any theorems.
        \item The proofs can either appear in the main paper or the supplemental material, but if they appear in the supplemental material, the authors are encouraged to provide a short proof sketch to provide intuition. 
        \item Inversely, any informal proof provided in the core of the paper should be complemented by formal proofs provided in appendix or supplemental material.
        \item Theorems and Lemmas that the proof relies upon should be properly referenced. 
    \end{itemize}

    \item {\bf Experimental result reproducibility}
    \item[] Question: Does the paper fully disclose all the information needed to reproduce the main experimental results of the paper to the extent that it affects the main claims and/or conclusions of the paper (regardless of whether the code and data are provided or not)?
    \item[] Answer: \answerYes{} % Replace by \answerYes{}, \answerNo{}, or \answerNA{}.
    \item[] Justification: Section 5 and Appendix B give the data-generating law, split sizes, scores, seed derivation, and uncertainty estimates. The artifact contains all executable code and raw results.
    \item[] Guidelines:
    \begin{itemize}
        \item The answer \answerNA{} means that the paper does not include experiments.
        \item If the paper includes experiments, a \answerNo{} answer to this question will not be perceived well by the reviewers: Making the paper reproducible is important, regardless of whether the code and data are provided or not.
        \item If the contribution is a dataset and\slash or model, the authors should describe the steps taken to make their results reproducible or verifiable. 
        \item Depending on the contribution, reproducibility can be accomplished in various ways. For example, if the contribution is a novel architecture, describing the architecture fully might suffice, or if the contribution is a specific model and empirical evaluation, it may be necessary to either make it possible for others to replicate the model with the same dataset, or provide access to the model. In general. releasing code and data is often one good way to accomplish this, but reproducibility can also be provided via detailed instructions for how to replicate the results, access to a hosted model (e.g., in the case of a large language model), releasing of a model checkpoint, or other means that are appropriate to the research performed.
        \item While NeurIPS does not require releasing code, the conference does require all submissions to provide some reasonable avenue for reproducibility, which may depend on the nature of the contribution. For example
        \begin{enumerate}
            \item If the contribution is primarily a new algorithm, the paper should make it clear how to reproduce that algorithm.
            \item If the contribution is primarily a new model architecture, the paper should describe the architecture clearly and fully.
            \item If the contribution is a new model (e.g., a large language model), then there should either be a way to access this model for reproducing the results or a way to reproduce the model (e.g., with an open-source dataset or instructions for how to construct the dataset).
            \item We recognize that reproducibility may be tricky in some cases, in which case authors are welcome to describe the particular way they provide for reproducibility. In the case of closed-source models, it may be that access to the model is limited in some way (e.g., to registered users), but it should be possible for other researchers to have some path to reproducing or verifying the results.
        \end{enumerate}
    \end{itemize}


\item {\bf Open access to data and code}
    \item[] Question: Does the paper provide open access to the data and code, with sufficient instructions to faithfully reproduce the main experimental results, as described in supplemental material?
    \item[] Answer: \answerYes{} % Replace by \answerYes{}, \answerNo{}, or \answerNA{}.
    \item[] Justification: The accompanying artifact contains code, generated raw data, pinned versions, and one-command reproduction instructions in README.md and reproduce.sh. It requires no accounts or external datasets.
    \item[] Guidelines:
    \begin{itemize}
        \item The answer \answerNA{} means that paper does not include experiments requiring code.
        \item Please see the NeurIPS code and data submission guidelines (\url{https://neurips.cc/public/guides/CodeSubmissionPolicy}) for more details.
        \item While we encourage the release of code and data, we understand that this might not be possible, so \answerNo{} is an acceptable answer. Papers cannot be rejected simply for not including code, unless this is central to the contribution (e.g., for a new open-source benchmark).
        \item The instructions should contain the exact command and environment needed to run to reproduce the results. See the NeurIPS code and data submission guidelines (\url{https://neurips.cc/public/guides/CodeSubmissionPolicy}) for more details.
        \item The authors should provide instructions on data access and preparation, including how to access the raw data, preprocessed data, intermediate data, and generated data, etc.
        \item The authors should provide scripts to reproduce all experimental results for the new proposed method and baselines. If only a subset of experiments are reproducible, they should state which ones are omitted from the script and why.
        \item At submission time, to preserve anonymity, the authors should release anonymized versions (if applicable).
        \item Providing as much information as possible in supplemental material (appended to the paper) is recommended, but including URLs to data and code is permitted.
    \end{itemize}


\item {\bf Experimental setting/details}
    \item[] Question: Does the paper specify all the training and test details (e.g., data splits, hyperparameters, how they were chosen, type of optimizer) necessary to understand the results?
    \item[] Answer: \answerYes{} % Replace by \answerYes{}, \answerNo{}, or \answerNA{}.
    \item[] Justification: Section 5 and Appendix B specify every score, concentration, smoothing setting, pilot estimate, and test distribution; there is no trained neural predictor.
    \item[] Guidelines:
    \begin{itemize}
        \item The answer \answerNA{} means that the paper does not include experiments.
        \item The experimental setting should be presented in the core of the paper to a level of detail that is necessary to appreciate the results and make sense of them.
        \item The full details can be provided either with the code, in appendix, or as supplemental material.
    \end{itemize}

\item {\bf Experiment statistical significance}
    \item[] Question: Does the paper report error bars suitably and correctly defined or other appropriate information about the statistical significance of the experiments?
    \item[] Answer: \answerYes{} % Replace by \answerYes{}, \answerNo{}, or \answerNA{}.
    \item[] Justification: Tables and figures report pointwise 95 percent intervals over independent repetition-level estimates; paired contrasts use paired repetitions. Appendix B explains the intervals and their limits.
    \item[] Guidelines:
    \begin{itemize}
        \item The answer \answerNA{} means that the paper does not include experiments.
        \item The authors should answer \answerYes{} if the results are accompanied by error bars, confidence intervals, or statistical significance tests, at least for the experiments that support the main claims of the paper.
        \item The factors of variability that the error bars are capturing should be clearly stated (for example, train/test split, initialization, random drawing of some parameter, or overall run with given experimental conditions).
        \item The method for calculating the error bars should be explained (closed form formula, call to a library function, bootstrap, etc.)
        \item The assumptions made should be given (e.g., Normally distributed errors).
        \item It should be clear whether the error bar is the standard deviation or the standard error of the mean.
        \item It is OK to report 1-sigma error bars, but one should state it. The authors should preferably report a 2-sigma error bar than state that they have a 96\% CI, if the hypothesis of Normality of errors is not verified.
        \item For asymmetric distributions, the authors should be careful not to show in tables or figures symmetric error bars that would yield results that are out of range (e.g., negative error rates).
        \item If error bars are reported in tables or plots, the authors should explain in the text how they were calculated and reference the corresponding figures or tables in the text.
    \end{itemize}

\item {\bf Experiments compute resources}
    \item[] Question: For each experiment, does the paper provide sufficient information on the computer resources (type of compute workers, memory, time of execution) needed to reproduce the experiments?
    \item[] Answer: \answerYes{} % Replace by \answerYes{}, \answerNo{}, or \answerNA{}.
    \item[] Justification: Section 5 reports CPU model, workers, elapsed time, and peak per-process memory. The artifact retains timing logs for the smoke, initial, final, and numerical-check runs.
    \item[] Guidelines:
    \begin{itemize}
        \item The answer \answerNA{} means that the paper does not include experiments.
        \item The paper should indicate the type of compute workers CPU or GPU, internal cluster, or cloud provider, including relevant memory and storage.
        \item The paper should provide the amount of compute required for each of the individual experimental runs as well as estimate the total compute. 
        \item The paper should disclose whether the full research project required more compute than the experiments reported in the paper (e.g., preliminary or failed experiments that didn't make it into the paper). 
    \end{itemize}
    
\item {\bf Code of ethics}
    \item[] Question: Does the research conducted in the paper conform, in every respect, with the NeurIPS Code of Ethics \url{https://neurips.cc/public/EthicsGuidelines}?
    \item[] Answer: \answerYes{} % Replace by \answerYes{}, \answerNo{}, or \answerNA{}.
    \item[] Justification: We reviewed the official NeurIPS Code of Ethics. This synthetic study uses no human data; Section 7 discusses misuse of marginal coverage and limitations.
    \item[] Guidelines:
    \begin{itemize}
        \item The answer \answerNA{} means that the authors have not reviewed the NeurIPS Code of Ethics.
        \item If the authors answer \answerNo, they should explain the special circumstances that require a deviation from the Code of Ethics.
        \item The authors should make sure to preserve anonymity (e.g., if there is a special consideration due to laws or regulations in their jurisdiction).
    \end{itemize}


\item {\bf Broader impacts}
    \item[] Question: Does the paper discuss both potential positive societal impacts and negative societal impacts of the work performed?
    \item[] Answer: \answerYes{} % Replace by \answerYes{}, \answerNo{}, or \answerNA{}.
    \item[] Justification: Section 7 describes potential benefits of support diagnostics and risks of false assurance, conservatism, and inappropriate fairness interpretations.
    \item[] Guidelines:
    \begin{itemize}
        \item The answer \answerNA{} means that there is no societal impact of the work performed.
        \item If the authors answer \answerNA{} or \answerNo, they should explain why their work has no societal impact or why the paper does not address societal impact.
        \item Examples of negative societal impacts include potential malicious or unintended uses (e.g., disinformation, generating fake profiles, surveillance), fairness considerations (e.g., deployment of technologies that could make decisions that unfairly impact specific groups), privacy considerations, and security considerations.
        \item The conference expects that many papers will be foundational research and not tied to particular applications, let alone deployments. However, if there is a direct path to any negative applications, the authors should point it out. For example, it is legitimate to point out that an improvement in the quality of generative models could be used to generate Deepfakes for disinformation. On the other hand, it is not needed to point out that a generic algorithm for optimizing neural networks could enable people to train models that generate Deepfakes faster.
        \item The authors should consider possible harms that could arise when the technology is being used as intended and functioning correctly, harms that could arise when the technology is being used as intended but gives incorrect results, and harms following from (intentional or unintentional) misuse of the technology.
        \item If there are negative societal impacts, the authors could also discuss possible mitigation strategies (e.g., gated release of models, providing defenses in addition to attacks, mechanisms for monitoring misuse, mechanisms to monitor how a system learns from feedback over time, improving the efficiency and accessibility of ML).
    \end{itemize}
    
\item {\bf Safeguards}
    \item[] Question: Does the paper describe safeguards that have been put in place for responsible release of data or models that have a high risk for misuse (e.g., pre-trained language models, image generators, or scraped datasets)?
    \item[] Answer: \answerNA{} % Replace by \answerYes{}, \answerNo{}, or \answerNA{}.
    \item[] Justification: No high-risk model, personal data, or scraped dataset is released. The assets are synthetic numerical observations and evaluation code.
    \item[] Guidelines:
    \begin{itemize}
        \item The answer \answerNA{} means that the paper poses no such risks.
        \item Released models that have a high risk for misuse or dual-use should be released with necessary safeguards to allow for controlled use of the model, for example by requiring that users adhere to usage guidelines or restrictions to access the model or implementing safety filters. 
        \item Datasets that have been scraped from the Internet could pose safety risks. The authors should describe how they avoided releasing unsafe images.
        \item We recognize that providing effective safeguards is challenging, and many papers do not require this, but we encourage authors to take this into account and make a best faith effort.
    \end{itemize}

\item {\bf Licenses for existing assets}
    \item[] Question: Are the creators or original owners of assets (e.g., code, data, models), used in the paper, properly credited and are the license and terms of use explicitly mentioned and properly respected?
    \item[] Answer: \answerNo{} % Replace by \answerYes{}, \answerNo{}, or \answerNA{}.
    \item[] Justification: All papers and package creators are credited and the artifact records package licenses and official-template provenance. We could not identify an explicit license in the official NeurIPS template archive, so no redistribution license is inferred for it.
    \item[] Guidelines:
    \begin{itemize}
        \item The answer \answerNA{} means that the paper does not use existing assets.
        \item The authors should cite the original paper that produced the code package or dataset.
        \item The authors should state which version of the asset is used and, if possible, include a URL.
        \item The name of the license (e.g., CC-BY 4.0) should be included for each asset.
        \item For scraped data from a particular source (e.g., website), the copyright and terms of service of that source should be provided.
        \item If assets are released, the license, copyright information, and terms of use in the package should be provided. For popular datasets, \url{paperswithcode.com/datasets} has curated licenses for some datasets. Their licensing guide can help determine the license of a dataset.
        \item For existing datasets that are re-packaged, both the original license and the license of the derived asset (if it has changed) should be provided.
        \item If this information is not available online, the authors are encouraged to reach out to the asset's creators.
    \end{itemize}

\item {\bf New assets}
    \item[] Question: Are new assets introduced in the paper well documented and is the documentation provided alongside the assets?
    \item[] Answer: \answerYes{} % Replace by \answerYes{}, \answerNo{}, or \answerNA{}.
    \item[] Justification: README.md, PROTOCOL.md, RESEARCH\_LOG.md, environment pins, and source docstrings document the new code and generated data, including intended use, limitations, and licenses.
    \item[] Guidelines:
    \begin{itemize}
        \item The answer \answerNA{} means that the paper does not release new assets.
        \item Researchers should communicate the details of the dataset\slash code\slash model as part of their submissions via structured templates. This includes details about training, license, limitations, etc. 
        \item The paper should discuss whether and how consent was obtained from people whose asset is used.
        \item At submission time, remember to anonymize your assets (if applicable). You can either create an anonymized URL or include an anonymized zip file.
    \end{itemize}

\item {\bf Crowdsourcing and research with human subjects}
    \item[] Question: For crowdsourcing experiments and research with human subjects, does the paper include the full text of instructions given to participants and screenshots, if applicable, as well as details about compensation (if any)? 
    \item[] Answer: \answerNA{} % Replace by \answerYes{}, \answerNo{}, or \answerNA{}.
    \item[] Justification: No crowdsourcing or human-subject experiment was performed.
    \item[] Guidelines:
    \begin{itemize}
        \item The answer \answerNA{} means that the paper does not involve crowdsourcing nor research with human subjects.
        \item Including this information in the supplemental material is fine, but if the main contribution of the paper involves human subjects, then as much detail as possible should be included in the main paper. 
        \item According to the NeurIPS Code of Ethics, workers involved in data collection, curation, or other labor should be paid at least the minimum wage in the country of the data collector. 
    \end{itemize}

\item {\bf Institutional review board (IRB) approvals or equivalent for research with human subjects}
    \item[] Question: Does the paper describe potential risks incurred by study participants, whether such risks were disclosed to the subjects, and whether Institutional Review Board (IRB) approvals (or an equivalent approval/review based on the requirements of your country or institution) were obtained?
    \item[] Answer: \answerNA{} % Replace by \answerYes{}, \answerNo{}, or \answerNA{}.
    \item[] Justification: The study uses only synthetic observations and has no human participants or personal data.
    \item[] Guidelines:
    \begin{itemize}
        \item The answer \answerNA{} means that the paper does not involve crowdsourcing nor research with human subjects.
        \item Depending on the country in which research is conducted, IRB approval (or equivalent) may be required for any human subjects research. If you obtained IRB approval, you should clearly state this in the paper. 
        \item We recognize that the procedures for this may vary significantly between institutions and locations, and we expect authors to adhere to the NeurIPS Code of Ethics and the guidelines for their institution. 
        \item For initial submissions, do not include any information that would break anonymity (if applicable), such as the institution conducting the review.
    \end{itemize}

\item {\bf Declaration of LLM usage}
    \item[] Question: Does the paper describe the usage of LLMs if it is an important, original, or non-standard component of the core methods in this research? Note that if the LLM is used only for writing, editing, or formatting purposes and does \emph{not} impact the core methodology, scientific rigor, or originality of the research, declaration is not required.
    %this research? 
    \item[] Answer: \answerYes{} % Replace by \answerYes{}, \answerNo{}, or \answerNA{}.
    \item[] Justification: Section 7 discloses autonomous LLM involvement in research design, implementation, analysis, and writing. LLM calls are not part of any experimental method, and no human review is claimed.
    \item[] Guidelines:
    \begin{itemize}
        \item The answer \answerNA{} means that the core method development in this research does not involve LLMs as any important, original, or non-standard components.
        \item Please refer to our LLM policy in the NeurIPS handbook for what should or should not be described.
    \end{itemize}

\end{enumerate}
\end{document}
